\SGAchapter{II}{整体消解的存在性}{Existence de résolutions globales}
\begin{center}L. Illusie\end{center}
\SourcePage{170}
\SGAsection{1}{整体化的一般判据}
\textbf{1.0。}本节给定一个\Term{site}{景} $\mathcal S$、一个\Term{flat-abelian-fibered-category}{平坦阿贝尔 $\mathcal S$-范畴} $\mathcal C$（第 I 讲 1.1），以及一个\Term{strictly-full-subcategory}{严格全}\Term{additive-category}{加性子范畴} $\mathcal C_0$（第 I 讲 1.1）。假定 $\mathcal C_0$ 在 $\mathcal C$ 中\Term{locally-quasi-liftable}{局部拟可提升}（第 I 讲 1.2），并且对\Term{epimorphism}{满态射}的\Term{kernel}{核}\Term{locally-quasi-stable}{局部拟稳定}（第 I 讲 1.2）。
\begin{SGAresult}{命题}{1.1}
设 $S\in\operatorname{ob}\mathcal S$，并假定：

(i) $\mathcal C_{0S}$ 在 $\mathcal C_S$ 中\Term{quasi-liftable}{拟可提升}（第 I 讲 1.2；这里把 $\mathcal C_S$ 看作点\Term{site}{景}上的\Term{fibered-category}{纤维范畴}）；

(ii) $\mathcal C_S$ 中每个为 $\mathcal C_0$-\Term{finite-type}{有限型}（第 I 讲 1.2）（或 $\mathcal C_0$-\Term{finite-presentation}{有限表现}，第 I 讲 2.8）的\Term{object}{对象}都是 $\mathcal C_{0S}$-\Term{finite-type}{有限型}的。

则：

a) $\mathcal C_{0S}$ 对\Term{epimorphism}{满态射}的\Term{kernel}{核}\Term{quasi-stable}{拟稳定}（第 I 讲 1.2；把 $\mathcal C_{0S}$ 看作点\Term{site}{景}上的\Term{fibered-category}{纤维范畴}）。

b) 设 $E\in\operatorname{ob}D^-(\mathcal C_S)$ 且 $n\in\mathbb Z$。若 $E$ 是 $n$-$\mathcal C_0$-\Term{pseudo-coherent}{伪凝聚}的（第 I 讲 2.3），则 $E$ 是 $n$-$\mathcal C_{0S}$-\Term{pseudo-coherent}{伪凝聚}的（在\Term{lie-bracket}{括号}中相应的假设下，是 $(n+1)$-$\mathcal C_{0S}$-\Term{pseudo-coherent}{伪凝聚}的）。换言之，
\[
D^-(\mathcal C_S)_{n\text{-}\mathcal C_0\text{-coh}}=D^-(\mathcal C_S)_{n\text{-}\mathcal C_{0S}\text{-coh}}
\]
（相应地，
\[
D^-(\mathcal C_S)_{n\text{-}\mathcal C_0\text{-coh}}
\subset D^-(\mathcal C_S)_{(n+1)\text{-}\mathcal C_{0S}\text{-coh}}
\subset D^-(\mathcal C_S)_{(n+1)\text{-}\mathcal C_0\text{-coh}},
\]
这些都是 $D(\mathcal C_S)$ 的\Term{strictly-full-subcategory}{严格全子范畴}的包含，右边的包含本来就是显然的）。

c) $E\in\operatorname{ob}D^-(\mathcal C_S)$ 相对于 $\mathcal C_0$ \Term{pseudo-coherent}{伪凝聚}（第 I 讲 2.3）的充要条件，是它相对于 $\mathcal C_{0S}$ \Term{pseudo-coherent}{伪凝聚}；即
\[
D^-(\mathcal C_S)_{\mathcal C_0\text{-coh}}=D^-(\mathcal C_S)_{\mathcal C_{0S}\text{-coh}}.
\]
此外，\Term{canonical}{典范}\Term{functor}{函子}
\[
K^-(\mathcal C_{0S})/K^{-,\varnothing}(\mathcal C_{0S})\longrightarrow D(\mathcal C_S)
\]
是\Term{fully-faithful}{全忠实}的，其\Term{essential-image}{本质像}为 $D^-(\mathcal C_S)_{\mathcal C_0\text{-coh}}$；这里 $K^{-,\varnothing}(\mathcal C_{0S})$ 表示 $K^-(\mathcal C_{0S})$ 中由\Term{acyclic-complex}{无上同调复形}组成的\Term{thick-subcategory}{厚子范畴}。
\SourcePage{171}
若 $\mathcal C_{0S}$ 在 $\mathcal C_S$ 中\Term{liftable}{可提升}（第 I 讲 1.2，即 $\mathcal C_{0S}$ 的\Term{object}{对象}都是\Term{projective-object}{投射对象}），则 $K^{-,\varnothing}(\mathcal C_{0S})$ 只含\Term{zero-object}{零对象}，因而\Term{functor}{函子}
\[
K^-(\mathcal C_{0S})\longrightarrow D(\mathcal C_S)
\]
是\Term{fully-faithful}{全忠实}的（其\Term{essential-image}{本质像}为 $D^-(\mathcal C_S)_{\mathcal C_0\text{-coh}}$）。
\end{SGAresult}
\begin{proof}
a) 设
\[
0\longrightarrow L\longrightarrow L^0\longrightarrow\cdots\longrightarrow L^m\longrightarrow0
\]
是 $\mathcal C_S$ 中的\Term{exact-sequence}{正合列}，且每个 $L^i\in\operatorname{ob}\mathcal C_{0S}$；须证 $L$ 为 $\mathcal C_{0S}$-\Term{finite-type}{有限型}。由第 I 讲 2.2，$L$ 为 $\mathcal C_0$-\Term{pseudo-coherent}{伪凝聚}，从而由第 I 讲 2.9 为 $\mathcal C_0$-\Term{finite-presentation}{有限表现}，于是结论由 (ii) 得出。

b) 对 $i$ 作\Term{descent}{下降}归纳，证明 $E$ 对 $i\geq n$（相应地 $i\geq n+1$）是 $i$-$\mathcal C_{0S}$-\Term{pseudo-coherent}{伪凝聚}的。可以先假设 $E$ 是严格 $(n+1)$-$\mathcal C_{0S}$-\Term{pseudo-coherent}{伪凝聚}的（相应地严格 $(n+2)$-$\mathcal C_{0S}$-\Term{pseudo-coherent}{伪凝聚}的，第 I 讲 2.1）；须证它是 $n$-$\mathcal C_{0S}$-\Term{pseudo-coherent}{伪凝聚}的（相应地 $(n+1)$-$\mathcal C_{0S}$-\Term{pseudo-coherent}{伪凝聚}的）。因此有\Term{exact-sequence}{正合列}
\[
0\longrightarrow E'\longrightarrow E\longrightarrow E''\longrightarrow0,
\]
其中 $E'$ 严格 $\mathcal C_{0S}$-\Term{perfect-complex}{完美}（第 I 讲 2.1），且 $E''{}^i=0$ 对 $i>n$（相应地 $i>n+1$）成立。用 $E''$ 替代 $E$，即可由第 I 讲 2.6 假定 $E^i=0$ 对 $i>n$（相应地 $i>n+1$）成立。于是由第 I 讲 2.10 a)，$H^n(E)$（相应地 $H^{n+1}(E)$）为 $\mathcal C_0$-\Term{finite-type}{有限型}（相应地 $\mathcal C_0$-\Term{finite-presentation}{有限表现}），由 (ii) 它是 $\mathcal C_{0S}$-\Term{finite-type}{有限型}。再应用第 I 讲 2.10 b)，便知 $E$ 是 $n$-$\mathcal C_{0S}$-\Term{pseudo-coherent}{伪凝聚}的（相应地 $(n+1)$-$\mathcal C_{0S}$-\Term{pseudo-coherent}{伪凝聚}的），证毕。（注意：引用第 I 讲 2.6 和 2.10 时，已隐含地使用了假设 (i) 和断言 a)。）

c) 第一断言由 b) 得到，第二断言由第 I 讲 2.7 得到；最后一个断言显然。
\end{proof}
\begin{SGAresult}{命题}{1.2}
假定 $\mathcal C_0$ 在 $\mathcal C$ 中\Term{locally-liftable}{局部可提升}（第 I 讲 1.2），且对\Term{epimorphism}{满态射}的\Term{kernel}{核}\Term{local-scheme}{局部}\Term{stable}{稳定}（第 I 讲 1.2），并作假设 1.1 (i) 及 (ii) 中\Term{lie-bracket}{括号}内的相应假设。
\SourcePage{172}
以 $\mathcal C_{1S}$ 表示 $\mathcal C_S$ 中由\Term{local-scheme}{局部}属于 $\mathcal C_0$ 的\Term{object}{对象}组成的\Term{strictly-full-subcategory}{严格全子范畴}。则：

a) $\mathcal C_{1S}$ 对\Term{epimorphism}{满态射}的\Term{kernel}{核}\Term{stable}{稳定}，并且在 $\mathcal C_S$ 中\Term{quasi-liftable}{拟可提升}（第 I 讲 1.2；作为点\Term{site}{景}上的\Term{fibered-category}{纤维范畴}）。

b) 设 $E\in\operatorname{ob}D(\mathcal C_S)$，$[m,n]$ 为 $\mathbb Z$ 中的区间。$E$ 的 $\mathcal C_0$-\Term{perfect-amplitude}{完美振幅}包含在 $[m,n]$ 中，当且仅当 $E$ 在 $D(\mathcal C_S)$ 中\Term{isomorphism}{同构}于某个\Term{complex}{复形} $F$，满足 $F^i=0$ 对 $i\notin[m,n]$ 成立，$F^i\in\operatorname{ob}\mathcal C_{0S}$ 对 $i>m$ 成立，且 $F^m\in\operatorname{ob}\mathcal C_{1S}$。

c) \Term{canonical}{典范}\Term{functor}{函子}
\[
K^b(\mathcal C_{1S})/K^{b,\varnothing}(\mathcal C_{1S})\longrightarrow D(\mathcal C_S)
\]
是\Term{fully-faithful}{全忠实}的，其\Term{essential-image}{本质像}为 $D(\mathcal C_S)$ 中由具有\Term{finite}{有限} $\mathcal C_0$-\Term{perfect-amplitude}{完美振幅}的\Term{complex}{复形}组成的\Term{full-subcategory}{全子范畴}。

d) 若 $\mathcal C_{1S}$ 在 $\mathcal C_S$ 中不仅\Term{quasi-liftable}{拟可提升}，而且\Term{liftable}{可提升}，则 b) 中“$E$ 在 $D(\mathcal C_S)$ 中\Term{isomorphism}{同构}于 $F$”可替换为“存在\Term{quasi-isomorphism}{拟同构} $F\to E$”，且 c) 中 $K^{b,\varnothing}(\mathcal C_{1S})$ 只含\Term{zero-object}{零对象}。
\end{SGAresult}
\begin{proof}
a) $\mathcal C_{1S}$ 对\Term{epimorphism}{满态射}的\Term{kernel}{核}\Term{stable}{稳定}是显然的。另一方面，$\mathcal C_{1S}$ 的\Term{object}{对象}是 $\mathcal C_0$-\Term{pseudo-coherent}{伪凝聚}的，由 1.1 c) 也是 $\mathcal C_{0S}$-\Term{pseudo-coherent}{伪凝聚}的；第二个断言便由第 I 讲 2.14 得到（其中取 $\mathcal S$ 为点\Term{site}{景}）。

b) 只须证明必要性。由 1.1 c)，已知 $E$ 是 $\mathcal C_{0S}$-\Term{pseudo-coherent}{伪凝聚}的。由于 $E$ 在\Term{degree}{次数} $>n$ 时\Term{acyclic-complex}{无上同调}，由第 I 讲 2.10，用 $D^-(\mathcal C_S)$ 中\Term{isomorphism}{同构}的\Term{complex}{复形}替代 $E$，即可假定 $E$ 严格 $(m+1)$-$\mathcal C_{0S}$-\Term{pseudo-coherent}{伪凝聚}，且 $E^i=0$ 对 $i>n$ 成立。又因 $E$ 在\Term{degree}{次数} $<m$ 时\Term{acyclic-complex}{无上同调}，将它替换为\Term{truncation}{截断}
\[
0\longrightarrow E^m/B^m\longrightarrow E^{m+1}\longrightarrow\cdots\longrightarrow E^n\longrightarrow0,
\]
还可假定 $E^i=0$ 对 $i<m$ 成立。于是由第 I 讲 4.16，$E^m$ \Term{local-scheme}{局部}为 $\mathcal C_0$ 的\Term{object}{对象}，因而为 $\mathcal C_{1S}$ 的\Term{object}{对象}，断言得证。
\SourcePage{173}
c) 结合 1.2 a)，本断言由 1.2 b) 和第 I 讲 4.12.2 得到。

d) $K^{b,\varnothing}(\mathcal C_{1S})$ 只含\Term{zero-object}{零对象}是显然的。另一断言由第 I 讲 4.6（在点\Term{site}{景}上应用）得到。
\end{proof}
\SGAsection{2}{在若干模范畴上的应用}
\textbf{2.0。记号。}设 $(\mathcal S,A)$ 为\Term{ringed-topos}{赋环拓扑斯}。对于 $U\in\operatorname{ob}\mathcal S$，以 $\operatorname{Mod}(U)$ 表示左 $A_U$-\Term{module}{模}的\Term{category}{范畴}，以 $\operatorname{Lib}(U)$（相应地 $\operatorname{Loclib}(U)$）表示其中由\Term{finite-type}{有限型}\Term{free-module}{自由} $A_U$-\Term{module}{模}（相应地\Term{local-scheme}{局部}为\Term{free-module}{自由模}的\Term{direct-summand}{直和因子}）组成的\Term{full-subcategory}{全子范畴}。
\textbf{2.1。\Term{lifting}{提升}\Term{lemma}{引理}。}
\begin{SGAresult}{\Term{lemma}{引理}}{2.1.1}
（参见 EGA $0_{\mathrm I}$ 5.2.3。）设 $(\mathcal S,A)$ 为\Term{ringed-topos}{赋环拓扑斯}，$S\in\operatorname{ob}\mathcal S$。假定 $S$ \Term{quasi-compact}{拟紧}，即每个\Term{covering-family}{覆盖族} $(S_i\to S)$ 都可由一个\Term{finite}{有限}\Term{covering-family}{覆盖族}\Term{refinement}{加细}。

a) 设左 $A_S$-\Term{module}{模} $F$ 是沿\Term{filtered-category}{滤过范畴} $\Lambda$ 的 $A_S$-\Term{module}{模} $F_\lambda$ 的\Term{inductive-limit}{归纳极限}，且
\[
u=\varinjlim u_\lambda:F\longrightarrow G
\]
为\Term{epimorphism}{满态射}，其中 $G$ 是\Term{finite-type}{有限型} $A_S$-\Term{module}{模}。则存在 $\lambda\in\Lambda$，使 $u_\lambda$ 为\Term{epimorphism}{满态射}。

b) 设左 $A_S$-\Term{module}{模} $F$ 由其\Term{section}{截面}生成（即对适当的指标集 $I$，存在\Term{epimorphism}{满态射} $A_S^{(I)}\to F$）。若 $F$ 为\Term{finite-type}{有限型}，则对某个 $p\in\mathbb N$ 存在\Term{epimorphism}{满态射} $A_S^p\to F$。
\end{SGAresult}
\begin{proof}
a) 的证明与所引文献中的证明类似，留给读者作练习；b) 立即由 a) 得到。
\end{proof}
\begin{SGAresult}{\Term{lemma}{引理}}{2.1.2}
设 $(\mathcal S,A)$ 为\Term{ringed-topos}{赋环拓扑斯}，$\mathcal C$ 为左 $A$-\Term{module}{模}的\Term{fibered-category}{纤维 $\mathcal S$-范畴}，$\mathcal C_0$ 为\Term{finite-type}{有限型}\Term{free-module}{自由} $A$-\Term{module}{模}组成的\Term{fibered-category}{纤维子范畴}（第 I 讲 1.3 d)），$S\in\operatorname{ob}\mathcal S$；$\mathcal C'_S$ 为 $\mathcal C_S$ 的一个包含 $\mathcal C_{0S}$ 且对\Term{kernel}{核}、\Term{cokernel}{余核}和\Term{extension}{扩张}\Term{stable}{稳定}的\Term{strictly-full-subcategory}{严格全}\Term{additive-category}{加性子范畴}。假定
\SourcePage{174}
\[
H^1(S,F)=0\qquad\text{对每个 }F\in\operatorname{ob}\mathcal C'_S.
\]
则：

a) $\mathcal C'_S$ 中每个\Term{exact-sequence}{正合列}
\[
0\longrightarrow E\longrightarrow F\longrightarrow G\longrightarrow0,
\]
只要 $\mathcal Hom(G,E)\in\operatorname{ob}\mathcal C'_S$ 且 $G$ \Term{local-scheme}{局部}为\Term{finite-type}{有限型}\Term{free-module}{自由模}的\Term{direct-summand}{直和因子}，就是\Term{split-exact-sequence}{分裂}的。$\mathcal C_{0S}$ 在 $\mathcal C'_S$ 中\Term{liftable}{可提升}（第 I 讲 1.2）。

b) 以 $\operatorname{Mod}(A(S))$ 表示 $A(S)$-\Term{module}{模}的\Term{category}{范畴}。\Term{functor}{函子}
\[
\Gamma_S:\mathcal C'_S\longrightarrow\operatorname{Mod}(A(S))
\]
是\Term{exact-functor}{正合}的。此外：

(i) 若 $\mathcal C'_S$ 的每个\Term{object}{对象} $F$ 都有整体\Term{finite-presentation}{有限表现}
\[
A_S^p\longrightarrow A_S^q\longrightarrow F\longrightarrow0,
\]
则 $\Gamma_S$ \Term{fully-faithful}{全忠实}，其\Term{essential-image}{本质像}为 $\operatorname{Mod}(A(S))$ 中由\Term{finite-presentation}{有限表现} $A(S)$-\Term{module}{模}组成的\Term{full-subcategory}{全子范畴}；

(ii) 若 $\mathcal C'_S$ 包含所有\Term{free-module}{自由} $A_S$-\Term{module}{模}，其每个\Term{object}{对象} $F$ 均有整体\Term{presentation}{表现}
\[
A_S^{(I)}\longrightarrow A_S^{(J)}\longrightarrow F\longrightarrow0,
\]
且 $\Gamma_S$ 与\Term{filtered-inductive-limit}{滤过归纳极限}交换，\footnote{当 $\mathcal S$ \Term{local-scheme}{局部}为\Term{finite-type}{有限型}时（SGA 4 VI 1.2），或当 $\mathcal S$ 为一个\Term{compact-space}{紧空间}上的\Term{set}{集合}\Term{sheaf}{层}的\Term{topos}{拓扑斯}且 $S$ 为\Term{final-object}{终对象}时（T.F. II 3.10.1），此假设成立。}则 $\Gamma_S:\mathcal C'_S\to\operatorname{Mod}(A(S))$ 为\Term{equivalence-of-categories}{范畴等价}。
\end{SGAresult}
\begin{proof}
a) 因 $G$ \Term{local-scheme}{局部}为\Term{finite-type}{有限型}\Term{free-module}{自由模}的\Term{direct-summand}{直和因子}，\Term{spectral-sequence}{谱序列}
\[
E_2^{pq}=H^p(S,\mathcal Ext^q(G,E))\Longrightarrow\operatorname{Ext}^*(G,E)
\]
\Term{degeneration}{退化}（第 I 讲 1.3.1），给出\Term{isomorphism}{同构}
\[
H^1(S,\mathcal Hom(G,E))\simeq\operatorname{Ext}^1(G,E),
\]
故有第一个断言。
\SourcePage{175}
第二断言是其特例：对 $G\in\operatorname{ob}\mathcal C_{0S}$，即 $G$ \Term{free-module}{自由}且为\Term{finite-type}{有限型}，以及 $E\in\operatorname{ob}\mathcal C'_S$，有 $\mathcal Hom(G,E)\in\operatorname{ob}\mathcal C'_S$。

b) 第一个断言显然。现于 (i) 或 (ii) 的情形证明 $\Gamma_S$ 的\Term{fully-faithful}{全忠实}性。设 $F,G\in\operatorname{ob}\mathcal C'_S$，并取 $F$ 的一个整体\Term{presentation}{表现}
\[
A_S^{(I)}\longrightarrow A_S^{(J)}\longrightarrow F\longrightarrow0.\tag{*}
\]
（在 (i) 中 $I,J$ \Term{finite}{有限}。）应用\Term{functor}{函子} $\Gamma_S$；它在 $\mathcal C'_S$ 上\Term{exact-functor}{正合}，且在 (ii) 中与\Term{filtered-inductive-limit}{滤过归纳极限}交换，故得到\Term{exact-sequence}{正合列}
\[
A(S)^{(I)}\longrightarrow A(S)^{(J)}\longrightarrow\Gamma_S(F)\longrightarrow0.\tag{**}
\]
另一方面，对 (*)（相应地 (**)）应用\Term{left-exact-functor}{左正合函子} $\operatorname{Hom}_{A_S}(\,\cdot\,,G)$（相应地 $\operatorname{Hom}_{A(S)}(\,\cdot\,,\Gamma_S(F))$），得下列各行\Term{exact-sequence}{正合}的\Term{commutative-diagram}{交换图}：
\[
\begin{tikzcd}[column sep=small]
0\ar[r]&\operatorname{Hom}_{A_S}(F,G)\ar[r]\ar[d]&\operatorname{Hom}_{A_S}(A_S^{(J)},G)\ar[r]\ar[d]&\operatorname{Hom}_{A_S}(A_S^{(I)},G)\ar[d]\\
0\ar[r]&\operatorname{Hom}_{A(S)}(F,G)\ar[r]&\operatorname{Hom}_{A(S)}(A(S)^{(J)},\Gamma_S(G))\ar[r]&\operatorname{Hom}_{A(S)}(A(S)^{(I)},\Gamma_S(G)).
\end{tikzcd}
\]
右边两条竖\Term{arrow}{箭头}显然是\Term{isomorphism}{同构}，故左边竖\Term{arrow}{箭头}也是\Term{isomorphism}{同构}；这证明了 $\Gamma_S$ 的\Term{fully-faithful}{全忠实}性。情形 (i) 或 (ii) 中\Term{essential-image}{本质像}的刻画是直接的，留给读者。
\end{proof}
\begin{SGAresult}{\Term{lemma}{引理}}{2.1.3}
设 $\mathcal A,\mathcal B$ 为\Term{abelian-category}{阿贝尔范畴}，$\phi:\mathcal A\to\mathcal B$ 为\Term{exact-functor}{正合函子}。

a) 对 $M\in\operatorname{ob}\mathcal A$，下列条件\Term{equivalence}{等价}：

(i) 对每个 $L\in\operatorname{ob}\mathcal A$，由 $\phi$ 诱导的\Term{canonical}{典范}\Term{arrow}{箭头}
\[
\operatorname{Ext}_{\mathcal A}^n(L,M)\longrightarrow\operatorname{Ext}_{\mathcal B}^n(\phi(L),\phi(M))
\]
对每个 $n$ 都是\Term{isomorphism}{同构}；

(ii) 对每个 $L\in\operatorname{ob}\mathcal A$，\Term{canonical}{典范}\Term{arrow}{箭头}
\SourcePage{176}
\[
\operatorname{Hom}_{\mathcal A}(L,M)\longrightarrow\operatorname{Hom}_{\mathcal B}(\phi(L),\phi(M))
\]
是\Term{isomorphism}{同构}，且对每个 $n\geq1$ 和每个 $x\in\operatorname{Ext}_{\mathcal B}^n(\phi(L),\phi(M))$，存在 $\mathcal A$ 的\Term{epimorphism}{满态射} $L'\to L$ 使 $x$ 消去。

b) 下列条件\Term{equivalence}{等价}：

(i) 对任意 $L,M\in\operatorname{ob}\mathcal A$ 及 $n\in\mathbb Z$，\Term{canonical}{典范}\Term{arrow}{箭头}
\[
\operatorname{Ext}_{\mathcal A}^n(L,M)\longrightarrow\operatorname{Ext}_{\mathcal B}^n(\phi(L),\phi(M))
\]
是\Term{isomorphism}{同构}；

(ii) 由 $\phi$ 诱导的\Term{canonical}{典范}\Term{functor}{函子}
\[
D^b(\mathcal A)\longrightarrow D^b(\mathcal B)
\]
是\Term{fully-faithful}{全忠实}的，其\Term{essential-image}{本质像}为 $D^b(\mathcal B)$ 中由满足每个 $H^i(F)$ 都在 $\phi$ 的\Term{essential-image}{本质像}中的\Term{complex}{复形} $F$ 组成的\Term{full-subcategory}{全子范畴}。
\end{SGAresult}
\begin{proof}
a) 为证 (i)$\Rightarrow$(ii)，只须证明：对 $n\geq1$、$L\in\operatorname{ob}\mathcal A$ 以及 $y\in\operatorname{Ext}_{\mathcal A}^n(L,M)$，存在\Term{epimorphism}{满态射} $L'\to L$ 使 $y$ 消去。可使用 Yoneda 的解释（[18] III 3.2）：$y$ 是一个\Term{exact-sequence}{正合列}
\[
0\longrightarrow M\longrightarrow E_{n-1}\longrightarrow\cdots\longrightarrow E_0\longrightarrow L\longrightarrow0
\]
的类；\Term{epimorphism}{满态射} $E_0\to L$ 使 $y$ 消去（[18] III 3.2.5 和 3.2.8）。再证 (ii)$\Rightarrow$(i)。简记
\[
F^n=\operatorname{Ext}_{\mathcal A}^n(\,\cdot\,,M),\qquad G^n=\operatorname{Ext}_{\mathcal B}^n(\phi(\,\cdot\,),\phi(M)).
\]
\Term{effaceable}{可消去}性假设意味着，对 $n\geq1$ 和 $L\in\operatorname{ob}\mathcal A$，有
\[
G^n(L)=\varinjlim\operatorname{Coker}\bigl(G^{n-1}(L_0)\longrightarrow G^{n-1}(L_1)\bigr),
\]
其中\Term{inductive-limit}{归纳极限}取遍所有\Term{exact-sequence}{正合列}
\SourcePage{177}
$0\to L_1\to L_0\to L\to0$。由上述论证，同样有
\[
F^n(L)=\varinjlim\operatorname{Coker}\bigl(F^{n-1}(L_0)\longrightarrow F^{n-1}(L_1)\bigr).
\]
因 $\phi=\phi^0:F^0\to G^0$ 为\Term{isomorphism}{同构}，对 $n$ 归纳，得由 $\phi$ 诱导的\Term{morphism}{态射} $\phi^n:F^n\to G^n$ 对每个 $n$ 都是\Term{isomorphism}{同构}。

b) (i)$\Leftrightarrow$(ii) 由直接的逐层\Term{resolution}{分解}得到，留给读者。
\end{proof}
\textbf{注记 2.1.3.1。}当然，a) 有一个\Term{duality}{对偶}的\Term{equivalence}{等价}条件，留给读者表述。另一方面，若 $\mathcal A,\mathcal B$ 有足够多\Term{injective-object}{内射对象}，则 b) 的条件 (i)、(ii) \Term{equivalence}{等价}于：

(iii) 对 $L\in\operatorname{ob}D^-(\mathcal A)$ 和 $M\in\operatorname{ob}D^+(\mathcal A)$，\Term{canonical}{典范}\Term{arrow}{箭头}
\[
\mathbf R\operatorname{Hom}(L,M)\longrightarrow\mathbf R\operatorname{Hom}(\phi(L),\phi(M))
\]
为\Term{isomorphism}{同构}。由\Term{spectral-sequence}{谱序列}论证立即可见这一点。
\textbf{2.2。\Term{noetherian-scheme}{诺特概形}与\Term{divisorial-scheme}{除子型概形}。}

\textbf{记号 2.2.1。}设 $S$ 为\Term{scheme}{概形}。以 $\operatorname{Qcoh}(S)$（相应地 $\operatorname{Coh}(S)$）表示\Term{quasi-coherent}{拟凝聚}（相应地\Term{coherent}{凝聚}）$\mathcal O_S$-\Term{module}{模}的（\Term{abelian-category}{阿贝尔}）\Term{category}{范畴}。设 $[m,n]$ 为 $\mathbb Z$ 的区间。称 $D(\operatorname{Qcoh}(S))$ 的\Term{object}{对象} $E$ 为 $n$-\Term{pseudo-coherent}{伪凝聚}的（相应地\Term{pseudo-coherent}{伪凝聚}、\Term{perfect-amplitude}{完美振幅}包含于 $[m,n]$、\Term{perfect-complex}{完美}），若它相对于 $S$ 的 Zariski \Term{site}{景}上由 $U\mapsto\operatorname{Loclib}(U)$ 定义的\Term{fibered-category}{纤维范畴}是如此（第 I 讲 2.3、4.7 及本讲 2.0）。以 $D^-(\operatorname{Qcoh}(S))_{\mathrm{coh}}$（相应地 $D(\operatorname{Qcoh}(S))_{\mathrm{parf}}$）表示 $D(\operatorname{Qcoh}(S))$ 中由\Term{pseudo-coherent}{伪凝聚}（相应地\Term{perfect-complex}{完美}）\Term{complex}{复形}组成的\Term{full-subcategory}{全}\Term{triangulated-category}{三角子范畴}。
\begin{SGAresult}{命题}{2.2.2}
\footnote{这是一个附带结果，后续讨论不再使用；建议读者初读时略过。}设 $S$ 为\Term{noetherian-scheme}{诺特概形}。\Term{canonical}{典范}\Term{functor}{函子}
\SourcePage{178}
\[
D^-(\operatorname{Coh}(S))\longrightarrow D(\operatorname{Qcoh}(S))
\]
是\Term{fully-faithful}{全忠实}的，其\Term{essential-image}{本质像}为 $D^-(\operatorname{Qcoh}(S))_{\mathrm{coh}}$。
\end{SGAresult}
\begin{proof}
应用 1.1，取 $\mathcal S$ 为 $S$ 的 Zariski \Term{site}{景}，取 $\mathcal C$（相应地 $\mathcal C_0$）为\Term{fibered-category}{纤维范畴} $U\mapsto\operatorname{Qcoh}(U)$（相应地 $\operatorname{Coh}(U)$）。因 $\mathcal O_S$ \Term{coherent}{凝聚}，$\mathcal C_0$ 的\Term{object}{对象}均\Term{pseudo-coherent}{伪凝聚}（第 I 讲 3.5）；故 $D(\operatorname{Qcoh}(U))$ 的\Term{object}{对象}（$U$ 为 $S$ 的\Term{open-subset}{开集}）在 2.2.1 的意义下\Term{pseudo-coherent}{伪凝聚}，当且仅当它相对于 $\mathcal C_0$ \Term{pseudo-coherent}{伪凝聚}（第 I 讲 2.14.1）。由 1.1 c)，只须证明 1.1 的条件 (i) 以及 (ii) 中不带\Term{lie-bracket}{括号}的条件成立。(ii) 显然，因为\Term{coherent}{凝聚}\Term{module}{模}的\Term{quasi-coherent}{拟凝聚}\Term{quotient-module}{商模}仍\Term{coherent}{凝聚}，且\Term{coherent}{凝聚}性是\Term{local-nature}{局部性质}。对 (i)，使用 2.1.1 a) 和如下事实：每个\Term{quasi-coherent}{拟凝聚} $\mathcal O_S$-\Term{module}{模}都是\Term{coherent}{凝聚} $\mathcal O_S$-\Term{module}{模}的\Term{filtered-inductive-limit}{滤过归纳极限}（EGA I 9.4.9）。
\end{proof}
承认独立于本小节的 3.5 与 3.7，容易\Term{pushout}{推出}：
\begin{SGAresult}{推论}{2.2.2.1}
设 $S$ 为\Term{noetherian-scheme}{诺特概形}（相应地\Term{finite-dimensional}{有限维}\Term{noetherian-scheme}{诺特概形}）。\Term{canonical}{典范}\Term{functor}{函子}
\[
D^b(\operatorname{Coh}(S))\longrightarrow D(S)
\qquad\bigl(\text{相应地 }D^-(\operatorname{Coh}(S))\longrightarrow D(S)\bigr)
\]
是\Term{fully-faithful}{全忠实}的，其\Term{essential-image}{本质像}为 $D^b(S)_{\mathrm{coh}}$（相应地 $D^-(S)_{\mathrm{coh}}$）（第 I 讲 2.15）。
\end{SGAresult}
\begin{SGAresult}{命题}{2.2.3}
（参见 EGA II 4.5.2、4.5.5。）设 $S$ 为\Term{quasi-compact}{拟紧}\Term{quasi-separated}{拟分离}\Term{scheme}{概形}，$(L_i)_{i\in I}$ 为 $S$ 上的一族\Term{invertible-sheaf}{可逆层}。对每个 $\mathcal O_S$-\Term{module}{模} $F$ 和 $(i,n)\in I\times\mathbb Z$，以 $F_{(i,n)}$ 表示 $F\otimes L_i^{\otimes n}$ 中由它在 $S$ 上的\Term{section}{截面}生成的 $\mathcal O_S$-\Term{submodule}{子模}。下列条件\Term{equivalence}{等价}：
\SourcePage{179}
(i) 对 $f\in\Gamma(S,L_i^{\otimes n})$、$i\in I$、$n\in\mathbb N$，各\Term{open-subset}{开集} $S_f$ 构成 $S$ 的\Term{topology}{拓扑}的一组基。

(i') 其中为\Term{affine-scheme}{仿射}的 $S_f$ \Term{covering}{覆盖} $S$。

(ii) 对每个\Term{finite-type}{有限型}\Term{quasi-coherent}{拟凝聚} $\mathcal O_S$-\Term{module}{模} $F$，存在一族非负整数 $(r_i)_{i\in I}$、一族正整数 $(n_i)_{i\in I}$，以及\Term{epimorphism}{满态射}
\[
\bigoplus_{i\in I}\mathcal O_S^{r_i}\otimes L_i^{\otimes(-n_i)}\longrightarrow F.
\]
(iii) 对每个 $F\in\operatorname{ob}\operatorname{Qcoh}(S)$，有
\[
F=\sum_{(i,n)\in I\times\mathbb N^+}F_{(i,n)}\otimes L_i^{\otimes(-n)}
\]
（这里是 $F$ 的\Term{subsheaf}{子层}的和，并非\Term{direct-sum}{直和}）。

(iii') 与 (iii) 相同，但只要求 $F$ 为 $\mathcal O_S$ 的\Term{quasi-coherent}{拟凝聚}\Term{ideal}{理想}。
\end{SGAresult}
\begin{proof}
这只是对所引文献的改述；为方便读者，简述如下。

(i)$\Rightarrow$(i')：对 $x\in S$，取包含 $x$ 的\Term{affine-scheme}{仿射}\Term{open-subset}{开集} $U$。若\Term{open-subset}{开集} $S_f$ 满足 $x\in S_f\subset U$，其中 $f\in\Gamma(S,L_i^{\otimes n})$、$n>0$，则 $S_f$ 必\Term{affine-scheme}{仿射}，由下述\Term{lemma}{引理}可知。
\end{proof}
\begin{SGAresult}{\Term{lemma}{引理}}{2.2.3.1}
（EGA II 5.5.8。）设 $X$ 为\Term{scheme}{概形}，$L$ 为\Term{invertible-sheaf}{可逆} $\mathcal O_X$-\Term{module}{模}，$f\in\Gamma(X,L)$。\Term{canonical}{典范}\Term{immersion}{浸入} $i:X_f\to X$ 是\Term{affine-morphism}{仿射态射}。
\end{SGAresult}
\begin{proof}
这是显然的：问题在 $X$ 上是\Term{local-scheme}{局部}的，故可假定 $X$ \Term{affine-scheme}{仿射}且 $L\simeq\mathcal O_X$，结论即得。
\end{proof}
\begin{proof}[命题 2.2.3 的证明（续）]
(i')$\Rightarrow$(ii)：取 $S$ 的一个\Term{finite}{有限}\Term{covering}{覆盖} $\mathfrak U$，由\Term{affine-scheme}{仿射}\Term{open-subset}{开集} $S_{f_{ij}}$ 组成，其中 $f_{ij}\in\Gamma(S,L_i^{\otimes m_{ij}})$；置 $S_i=\bigcup_jS_{f_{ij}}$。将 $f_{ij}$ 换成适当的幂，可假定同一 $i$ 对应的 $m_{ij}$ 都等于一个整数 $m_i$。
\SourcePage{180}
\Term{sheaf}{层} $F|S_{f_{ij}}$ 由它在 $S_{f_{ij}}$ 上的\Term{finite}{有限}个\Term{section}{截面} $h_{ijk}$ 生成。存在整数 $q_i\geq0$，使 $h_{ijk}\otimes f_{ij}^{\otimes q_i}$ 延拓为 $F\otimes L_i^{\otimes q_im_i}$ 在 $S$ 上的\Term{section}{截面}（EGA I 9.3.1、IV 1.7.5）。换言之，$F\otimes L_i^{\otimes q_im_i}$ 在 $S_i$ 上由它在 $S$ 上的\Term{finite}{有限}个\Term{section}{截面}生成。仿照 EGA II 4.5.5 的证明，可知存在整数 $n_i$（可取为正），使对所有 $n\geq n_i$，$F\otimes L_i^{\otimes n}$ 在 $S_i$ 上由它在 $S$ 上的\Term{finite}{有限}个\Term{section}{截面}生成。特别地，存在整数 $r_i\geq0$ 和\Term{morphism}{态射}
\[
\mathcal O_S^{r_i}\otimes L_i^{\otimes(-n_i)}\longrightarrow F,
\]
在 $S_i$ 上诱导\Term{epimorphism}{满态射}。于是立即得到 (i')$\Rightarrow$(ii)。

(ii)$\Rightarrow$(iii)：因 $S$ 上每个\Term{quasi-coherent}{拟凝聚}\Term{sheaf}{层}都是其\Term{finite-type}{有限型}\Term{quasi-coherent}{拟凝聚}\Term{subsheaf}{子层}的\Term{inductive-limit}{归纳极限}（EGA I 9.4.9、IV 1.7.7），故只须对\Term{finite-type}{有限型} $F$ 证明 (iii)。此时存在\Term{morphism}{态射}
\[
u_i:\mathcal O_S^{r_i}\otimes L_i^{\otimes(-n_i)}\longrightarrow F\quad(i\in I),
\]
使 $\sum u_i$ 为\Term{epimorphism}{满态射}。又
\[
\begin{aligned}
\operatorname{Im}(u_i)&=\operatorname{Im}\bigl(u_i\otimes\operatorname{Id}(L_i^{\otimes n_i})\otimes\operatorname{Id}(L_i^{\otimes(-n_i)})\bigr)\\
&=\operatorname{Im}\bigl(u_i\otimes\operatorname{Id}(L_i^{\otimes n_i})\bigr)\otimes L_i^{\otimes(-n_i)}.
\end{aligned}
\]
其中 $\operatorname{Im}(u_i\otimes\operatorname{Id}(L_i^{\otimes n_i}))$ 为 $F_{(i,n_i)}$ 的\Term{subsheaf}{子层}，故断言成立。

(iii)$\Rightarrow$(iii') 显然。

(iii')$\Rightarrow$(i)：设 $x\in S$，$U$ 为包含 $x$ 的\Term{open-subset}{开集}，$F$ 为定义 $S-U$ 的\Term{quasi-coherent}{拟凝聚}\Term{ideal}{理想}。
\SourcePage{181}
由 (iii')，存在 $(i,n)\in I\times\mathbb N^+$ 及 $F\otimes L_i^{\otimes n}$ 在 $S$ 上的\Term{section}{截面} $f$，使 $f(x)\ne0$。于是 $x\in S_f\subset U$，证明了 (i)。
\end{proof}
\begin{SGAresult}{定义}{2.2.4}
设 $S$ 为\Term{quasi-compact}{拟紧}\Term{quasi-separated}{拟分离}\Term{scheme}{概形}，$(L_i)_{i\in I}$ 为 $S$ 上的一族\Term{invertible-sheaf}{可逆层}。若 2.2.3 的\Term{equivalence}{等价}条件成立，则称各 $L_i$ 构成一个\Term{ample-family}{丰沛族}。\footnote{M. Borelli 和 S. L. Kleiman 以 (i') 的\Term{form}{形式}引入了这一概念；“\Term{ample-family}{丰沛族}”这一名称由撰稿人提出。}（特别地，只含一个\Term{sheaf}{层} $L$ 的族\Term{ample-family}{丰沛}，当且仅当 $L$ 在通常意义下\Term{ample-sheaf}{丰沛}，EGA II 4.5.3。）
\end{SGAresult}
\begin{SGAresult}{定义}{2.2.5}
称\Term{scheme}{概形} $S$ 为\Term{divisorial-scheme}{除子型概形}（divisoriel），若 $S$ \Term{quasi-compact}{拟紧}\Term{quasi-separated}{拟分离}，且存在\Term{invertible-sheaf}{可逆} $\mathcal O_S$-\Term{module}{模}的\Term{ample-family}{丰沛族}（2.2.4）。
\end{SGAresult}
\begin{SGAresult}{\Term{lemma}{引理}}{2.2.6}
设 $X$ 为\Term{locally-noetherian-scheme}{局部诺特概形}，$U$ 为 $X$ 的\Term{open-subset}{开集}，且包含 $U\to X$ 为\Term{affine-morphism}{仿射态射}（例如 $U$ \Term{affine-scheme}{仿射}且 $X$ \Term{separated}{分离}）。则 $X-U$ 纯\Term{codimension}{余维}为 $0$ 或 $1$（即对 $X-U$ 的每个\Term{maximal-point}{极大点} $x$，$\dim\mathcal O_{X,x}=0$ 或 $1$）；所以，若 $U$ 包含 $X$ 的\Term{maximal-point}{极大点}，则它纯\Term{codimension}{余维}为 $1$。
\end{SGAresult}
\begin{proof}
置 $Y=X-U$。设 $x$ 为 $Y$ 的\Term{maximal-point}{极大点}，作\Term{base-change}{基变换} $X'=\operatorname{Spec}(\mathcal O_{X,x})\to X$，得到\Term{cartesian-diagram}{笛卡尔图}
\[
\begin{tikzcd}
U'\ar[r]\ar[d]&X'\ar[d]&Y'\ar[l]\ar[d]\\
U\ar[r]&X&Y\ar[l].
\end{tikzcd}
\]
其中 $Y'=X'-U'$。因为 $x$ 是 $Y$ 的\Term{maximal-point}{极大点}，$Y'$ 正是 $X'$ 的\Term{closed-point}{闭点}。另一方面，$X'$ \Term{affine-scheme}{仿射}，故 $U'$ 为 $X'$ 的\Term{affine-scheme}{仿射}\Term{open-subset}{开集}。证明 $X'$ 的\Term{dimension}{维数} $n$ 等于 $0$ 或 $1$。事实上，已知 $H_{Y'}^n(X',\mathcal O_{X'})\ne0$（SGA 2 V 5.7）；而若 $n\geq2$，则
\[
H_{Y'}^n(X',\mathcal O_{X'})\simeq H^{n-1}(U',\mathcal O_{U'}),\qquad H^{n-1}(U',\mathcal O_{U'})=0,
\]
\SourcePage{182}
因为 $U'$ \Term{affine-scheme}{仿射}且 $n-1\geq1$，矛盾。
\end{proof}
\textbf{注记 2.2.6.1。}若去掉 $X$ \Term{separated}{分离}的假设，则对 $X$ 的\Term{affine-scheme}{仿射}\Term{open-subset}{开集} $U$，2.2.6 不再成立；原点加倍的\Term{affine-space}{仿射平面}即为反例（EGA I 5.5.11）。
\begin{SGAresult}{命题}{2.2.7}
设 $S$ 为\Term{noetherian-scheme}{诺特}、\Term{separated}{分离}、\Term{locally-factorial}{局部析因}的\Term{scheme}{概形}（即其\Term{local-ring}{局部环}均为\Term{factorial-ring}{析因环}）。则 $S$ 是\Term{divisorial-scheme}{除子型概形}（2.2.5）。
\end{SGAresult}
特别地（EGA IV 21.11.1）：
\begin{SGAresult}{推论}{2.2.7.1}
每个\Term{noetherian-scheme}{诺特}\Term{regular-scheme}{正则}\Term{separated}{分离}\Term{scheme}{概形}都是\Term{divisorial-scheme}{除子型概形}。
\end{SGAresult}
\begin{proof}[2.2.7 的证明]
可假设 $S$ \Term{integral-scheme}{整}。取 $S$ 的\Term{finite}{有限}\Term{covering}{覆盖} $(S_i)_{i\in I}$，各 $S_i$ 为\Term{nonempty}{非空}\Term{affine-scheme}{仿射}\Term{open-subset}{开集}。由 2.2.6，$S-S_i$ 纯\Term{codimension}{余维}为 $1$，从而由 EGA IV 21.7.4，它是一个\Term{effective-divisor}{有效除子}的\Term{support}{支集}。换言之，对每个 $i\in I$，存在\Term{invertible-sheaf}{可逆层} $L_i$ 及\Term{section}{截面} $f_i\in\Gamma(S,L_i)$，使 $S_i=S_{f_i}$。由 2.2.3 (i')，族 $(L_i)$ \Term{ample-family}{丰沛}（2.2.4）。证明完毕。
\end{proof}
\textbf{注记 2.2.7.2。}与 2.2.6 相同，$S$ \Term{separated}{分离}的假设不可少。仍取 $S$ 为原点加倍的\Term{affine-space}{仿射平面}。回顾其构造：将两个\Term{affine-space}{仿射平面} $\mathbb E_k^2$（$k$ 为域）沿原点 $O$ 的补集 $U$ 以\Term{identity-map}{恒等映射}\Term{gluing}{粘合}，得到 $S$；因此有\Term{canonical-projection}{典范投影}
\[
p:S\longrightarrow\mathbb E_k^2,
\]
它在 $U$ 上诱导\Term{isomorphism}{同构}，而 $p^{-1}(O)$ 由两点 $O_1,O_2$ 组成。$S$ 上的\Term{invertible-sheaf}{可逆层}由 $\mathbb E_k^2$ 上两个\Term{invertible-sheaf}{可逆层}沿 $U$ 上的\Term{isomorphism}{同构}\Term{gluing}{粘合}而来。因 $\mathbb E_k^2$ 上每个\Term{invertible-sheaf}{可逆层}都\Term{trivial-bundle}{平凡}，$S$ 上的\Term{invertible-sheaf}{可逆层}由某个 $g\in\Gamma(U,\mathcal O_U^*)$ 定义。但 $g$ 必为常数：由于 $O$ 在平面中的\Term{depth}{深度}为 $2$，\Term{restriction-map}{限制映射}
\[
\Gamma(\mathbb E_k^2,\mathcal O_{\mathbb E_k^2})\longrightarrow\Gamma(U,\mathcal O_U)
\]
为\Term{isomorphism}{同构}。故 $S$ 上每个\Term{invertible-sheaf}{可逆层}都\Term{trivial-bundle}{平凡}。
\SourcePage{183}
另一方面，显然 $p$ 诱导\Term{isomorphism}{同构}
\[
p^*:\Gamma(\mathbb E_k^2,\mathcal O_{\mathbb E_k^2})\longrightarrow\Gamma(S,\mathcal O_S).
\]
因此，$\mathcal O_S$ 的每个在 $O_1$ 处为零的\Term{global-section}{整体截面}在 $O_2$ 处也为零，所以对 $f\in\Gamma(S,\mathcal O_S)$，各 $S_f$ 不构成 $S$ 的\Term{topology}{拓扑}的一组基。
\begin{SGAresult}{\Term{lemma}{引理}}{2.2.8}
设 $S$ 为\Term{divisorial-scheme}{除子型概形}（2.2.5），$(L_i)_{i\in I}$ 为\Term{invertible-sheaf}{可逆} $\mathcal O_S$-\Term{module}{模}的\Term{ample-family}{丰沛族}。以 $\mathcal L(S)$ 表示 $\operatorname{Qcoh}(S)$（2.2.1）中由 $L_i^{\otimes(-n)}$、$(i,n)\in I\times\mathbb N^+$ 生成的\Term{strictly-full-subcategory}{严格全}\Term{additive-category}{加性子范畴}。设 $E\in\operatorname{ob}D(\operatorname{Qcoh}(S))$。

a) $\mathcal L(S)$ 及 $\operatorname{Loclib}(S)$（2.0）在 $\operatorname{Qcoh}(S)$ 中\Term{quasi-liftable}{拟可提升}（第 I 讲 1.2），且对\Term{epimorphism}{满态射}的\Term{kernel}{核}\Term{quasi-stable}{拟稳定}（第 I 讲 1.2）。

b) 设 $n\in\mathbb Z$。$E$ 为 $n$-\Term{pseudo-coherent}{伪凝聚}（相应地\Term{pseudo-coherent}{伪凝聚}，2.2.1）的充要条件，是它相对于 $\mathcal L(S)$ 或 $\operatorname{Loclib}(S)$ 为 $n$-\Term{pseudo-coherent}{伪凝聚}（相应地\Term{pseudo-coherent}{伪凝聚}，第 I 讲 2.3）。\Term{canonical}{典范}\Term{functor}{函子}
\[
K^-(\mathcal L(S))/K^{-,\varnothing}(\mathcal L(S))\longrightarrow D(\operatorname{Qcoh}(S))
\]
及
\[
K^-(\operatorname{Loclib}(S))/K^{-,\varnothing}(\operatorname{Loclib}(S))\longrightarrow D(\operatorname{Qcoh}(S))
\]
均\Term{fully-faithful}{全忠实}，\Term{essential-image}{本质像}为 $D^-(\operatorname{Qcoh}(S))_{\mathrm{coh}}$（2.2.1）。

c) 设 $[m,n]\subset\mathbb Z$ 为区间。$E$ 的\Term{perfect-amplitude}{完美振幅}包含在 $[m,n]$（2.2.1）中的充要条件，是 $E$ 在 $D(\operatorname{Qcoh}(S))$ 中\Term{isomorphism}{同构}于一个\Term{complex}{复形} $F$，使 $F^i=0$ 对 $i\notin[m,n]$ 成立，$F^i\in\operatorname{ob}\mathcal L(S)$ 对 $i>m$ 成立，且 $F^m$ \Term{locally-free}{局部自由}且为\Term{finite-type}{有限型}。\Term{canonical}{典范}\Term{functor}{函子}
\[
K^b(\operatorname{Loclib}(S))/K^{b,\varnothing}(\operatorname{Loclib}(S))\longrightarrow D(\operatorname{Qcoh}(S))
\]
\SourcePage{184}
是\Term{fully-faithful}{全忠实}的，\Term{essential-image}{本质像}为 $D(\operatorname{Qcoh}(S))_{\mathrm{parf}}$（2.2.1）。
\end{SGAresult}
\begin{proof}
设 $F\to G$ 为 $\operatorname{Qcoh}(S)$ 中的\Term{epimorphism}{满态射}，且 $G$ 为\Term{finite-type}{有限型}。因 $F$ 为其\Term{finite-type}{有限型}\Term{quasi-coherent}{拟凝聚}\Term{subsheaf}{子层}的\Term{inductive-limit}{归纳极限}（EGA I 9.4.9），由 2.1.1 a)，存在 $F$ 的\Term{finite-type}{有限型}\Term{quasi-coherent}{拟凝聚}\Term{subsheaf}{子层} $F'$，使复合 $F'\to F\to G$ 为\Term{epimorphism}{满态射}。由 2.2.3 (ii)，存在\Term{epimorphism}{满态射} $E\to F'$，其中 $E\in\operatorname{ob}\mathcal L(S)$；故 $E\to F'\to F$ 与 $F\to G$ 的复合为\Term{epimorphism}{满态射}。这说明\Term{finite-type}{有限型}\Term{quasi-coherent}{拟凝聚} $\mathcal O_S$-\Term{module}{模}的\Term{category}{范畴}在 $\operatorname{Qcoh}(S)$ 中\Term{quasi-liftable}{拟可提升}，从而 $\operatorname{Loclib}(S)$ 或 $\mathcal L(S)$ 更是如此。这就证明了 a) 的第一断言。

为证第二断言，应用 1.1，取 $\mathcal S$ 为 $S$ 的 Zariski \Term{site}{景}，$\mathcal C$（相应地 $\mathcal C_0$）为\Term{fibered-category}{纤维范畴} $U\mapsto\operatorname{Qcoh}(U)$（相应地 $\mathcal L(U)$）；这里 $U$ 遍历 $S$ 的\Term{open-subset}{开集}，$\mathcal L(U)$ 为由各 $L_i^{\otimes(-n)}|U$ 生成的\Term{strictly-full-subcategory}{严格全}\Term{additive-category}{加性子范畴}。因 $\mathcal L$ 的每个\Term{object}{对象}都\Term{pseudo-coherent}{伪凝聚}，\Term{fibered-category}{纤维范畴} $D(\operatorname{Qcoh})$ 中 $n$-\Term{pseudo-coherent}{伪凝聚}（相应地\Term{pseudo-coherent}{伪凝聚}）的概念（2.2.1）与相对于 $\mathcal L$ 的同一概念一致（第 I 讲 2.14）。由 2.2.3 (ii)，1.1 的条件 (ii) 中不带\Term{lie-bracket}{括号}的版本成立；于是由 1.1 a) 得到 a) 中关于 $\mathcal L(S)$ 的第二部分；关于 $\operatorname{Loclib}(S)$ 的断言同理。

b) 立即由 1.1 b)、c) 得到；这里须注意，若 $E\in\operatorname{ob}D(\operatorname{Qcoh}(S))$ 为 $n$-\Term{pseudo-coherent}{伪凝聚}，则由于 $S$ \Term{quasi-compact}{拟紧}，$E$ 是 $D^-(\operatorname{Qcoh}(S))$ 的\Term{object}{对象}。\Term{fibered-category}{纤维范畴} $\operatorname{Qcoh}$ 的\Term{object}{对象}\Term{local-scheme}{局部}属于 $\mathcal L$，当且仅当它\Term{locally-free}{局部自由}且为\Term{finite-type}{有限型}；所以 $\mathcal L$ 在 $\operatorname{Qcoh}$ 中\Term{locally-liftable}{局部可提升}，且对\Term{epimorphism}{满态射}的\Term{kernel}{核}\Term{local-scheme}{局部}\Term{stable}{稳定}。因此，取 $\mathcal C=\operatorname{Qcoh}$、$\mathcal C_0=\mathcal L$、$\mathcal C_{1S}=\operatorname{Loclib}(S)$，由 1.2 得到 c)。证明完毕。
\end{proof}
\begin{SGAresult}{命题}{2.2.9}
假设和记号与 2.2.8 相同，另假定 $S$ \Term{separated}{分离}或\Term{noetherian-scheme}{诺特}。

a) \Term{canonical}{典范}\Term{functor}{函子}
\[
K^{-,b}(\mathcal L(S))/K^{-,\varnothing}(\mathcal L(S))\longrightarrow D(S)
\]
及
\[
K^{-,b}(\operatorname{Loclib}(S))/K^{-,\varnothing}(\operatorname{Loclib}(S))\longrightarrow D(S)
\]
\SourcePage{185}
均\Term{fully-faithful}{全忠实}，\Term{essential-image}{本质像}为 $D^b(S)_{\mathrm{coh}}$（第 I 讲 2.15）。此外，若 $S$ 在 3.7 的意义下具有“\Term{finite}{有限}\Term{cohomological-dimension}{上同调维数}”\footnote{第 3 节独立于本讲其余部分。}（例如 $S$ 的底\Term{space}{空间}\Term{noetherian-space}{诺特}且具有\Term{finite}{有限} \Term{krull-dimension}{Krull 维数}），则\Term{functor}{函子}
\[
K^-(\mathcal L(S))/K^{-,\varnothing}(\mathcal L(S))\longrightarrow
K^-(\operatorname{Loclib}(S))/K^{-,\varnothing}(\operatorname{Loclib}(S))\longrightarrow D^-(S)_{\mathrm{coh}}
\]
均为\Term{equivalence-of-categories}{范畴等价}。

b) \Term{canonical}{典范}\Term{functor}{函子}
\[
K^b(\operatorname{Loclib}(S))/K^{b,\varnothing}(\operatorname{Loclib}(S))\longrightarrow D(S)
\]
是\Term{fully-faithful}{全忠实}的，其\Term{essential-image}{本质像}为 $D(S)_{\mathrm{parf}}$（第 I 讲 4.9）。$E\in\operatorname{ob}D(S)$ 的\Term{perfect-amplitude}{完美振幅}包含于 $[m,n]$（第 I 讲 4.8）中的充要条件，是 $E$ 在 $D(S)$ 中\Term{isomorphism}{同构}于一个\Term{complex}{复形} $F$，使 $F^i=0$ 对 $i\notin[m,n]$ 成立，$F^i\in\operatorname{ob}\mathcal L(S)$ 对 $i>m$ 成立，且 $F^m$ \Term{locally-free}{局部自由}且为\Term{finite-type}{有限型}。
\end{SGAresult}
\begin{proof}
结合 3.5 和 3.7（见上脚注），这是 2.2.8 的直接推论。唯一须说明的一点（它也显然由 3.5、相应地 3.7 得到）是：对 $E\in\operatorname{ob}D^b(\operatorname{Qcoh}(S))$（相应地，在 $S$ 有 3.7 意义下\Term{finite}{有限}\Term{cohomological-dimension}{上同调维数}时，对 $E\in\operatorname{ob}D(\operatorname{Qcoh}(S))$），把 $E$ 看作 $D(S)$ 的\Term{object}{对象}后，2.2.1 中的\Term{pseudo-coherent}{伪凝聚}性和\Term{perfect-complex}{完美}性与通常的概念（第 I 讲 2.15、4.9）一致。
\end{proof}
\textbf{注记 2.2.10。}a) 若 $S$ 为环 $A$ 的\Term{affine-scheme}{仿射概形}，则只含 $\mathcal O_S$ 的族\Term{ample-family}{丰沛}。对此族有 $\mathcal L(S)=\operatorname{Lib}(S)$。此外，显然 $K^{-,\varnothing}(\operatorname{Lib}(S))$ 和 $K^{-,\varnothing}(\operatorname{Loclib}(S))$ 都只含\Term{zero-object}{零对象}。因此由 2.2.9 得知，若以 $\operatorname{Mod}(A)$ 表示 $A$-\Term{module}{模}的\Term{category}{范畴}，且 $D(A)=D(\operatorname{Mod}(A))$，则\Term{functor}{函子} $\mathbf R\Gamma_S:D^+(S)\to D^+(A)$ 诱导\Term{fully-faithful}{全忠实}\Term{functor}{函子}
\[
D^b(S)_{\mathrm{coh}}\longrightarrow D(A)\qquad
(\text{相应地 }D(S)_{\mathrm{parf}}\longrightarrow D(A)),
\]
其\Term{essential-image}{本质像}为
\SourcePage{186}
\Term{triangulated-category}{三角子范畴} $D^b(A)_{\mathrm{coh}}$（相应地 $D(A)_{\mathrm{parf}}$），其\Term{object}{对象}为\Term{cohomology}{上同调}\Term{bounded}{有界}且相对于\Term{finite-type}{有限型}\Term{free-module}{自由}（相应地\Term{projective-module}{投射}）\Term{module}{模}的\Term{subcategory}{子范畴} $\operatorname{Lib}(A)$（相应地 $\operatorname{Loclib}(A)$）\Term{pseudo-coherent}{伪凝聚}（相应地\Term{perfect-complex}{完美}）的\Term{complex}{复形}。有由\Term{equivalence-of-categories}{范畴等价}组成的\Term{commutative-diagram}{交换图}
\[
\begin{tikzcd}
K^{-,b}(\operatorname{Lib}(S))\ar[r]\ar[d,"\Gamma_S"]&K^{-,b}(\operatorname{Loclib}(S))\ar[r]\ar[d,"\Gamma_S"]&D^b(S)_{\mathrm{coh}}\ar[d,"\mathbf R\Gamma_S"]\\
K^{-,b}(\operatorname{Lib}(A))\ar[r]&K^{-,b}(\operatorname{Loclib}(A))\ar[r]&D^b(A)_{\mathrm{coh}},
\end{tikzcd}
\]
\[
\begin{tikzcd}
K^b(\operatorname{Loclib}(S))\ar[r]\ar[d,"\Gamma_S"]&D(S)_{\mathrm{parf}}\ar[d,"\mathbf R\Gamma_S"]\\
K^b(\operatorname{Loclib}(A))\ar[r]&D(A)_{\mathrm{parf}}.
\end{tikzcd}
\]
（若 $S$ 在 3.7 的意义下有\Term{finite}{有限}\Term{cohomological-dimension}{上同调维数}，则上图中可将 $b$ 替换为 $-$。）

b) 设 $S$ 为\Term{scheme}{概形}。由 a)，一个\Term{pseudo-coherent}{伪凝聚}且\Term{cohomology}{上同调}\Term{bounded}{有界}的 $\mathcal O_S$-\Term{module}{模}\Term{complex}{复形}\Term{local-scheme}{局部}\Term{isomorphism}{同构}于（在 $D(S)$ 中）一个严格\Term{pseudo-coherent}{伪凝聚}\Term{complex}{复形}；这先验上远非显然（参见第 I 讲 2.3.1）。

\textbf{2.3。\Term{soft-sheaf}{软层}。}

\textbf{2.3.1。}回顾 [T.F.] 中的一些定义和结果。设 $X$ 为\Term{paracompact-space}{仿紧空间}。称 $X$ 上的\Term{set}{集合}\Term{sheaf}{层} $F$ 是\Term{soft-sheaf}{软层}，若对 $X$ 的每个\Term{closed-subset}{闭集} $Y$，$F$ 在 $Y$ 上的每个\Term{section}{截面}均可延拓为在 $X$ 上的\Term{section}{截面}。\Term{soft-sheaf}{软}\Term{abelian-sheaf}{阿贝尔层}是\Term{acyclic-sheaf}{无上同调}的，即若 $F$ \Term{soft-sheaf}{软}，则 $H^i(X,F)=0$ 对 $i>0$ 成立（T.F. II 4.4.3）。

设 $A$ 为 $X$ 上的\Term{sheaf-of-rings}{环层}。若 $F$ 为\Term{soft-sheaf}{软} $A$-\Term{module}{模}，则 $F$ 由其\Term{section}{截面}生成（参见 2.1.1）：由 $I=H^0(X,F)$ 定义的\Term{homomorphism}{同态} $A_X^{(I)}\to F$ 显然是\Term{epimorphism}{满态射}。
\SourcePage{187}
另一方面，若 $A$ \Term{soft-sheaf}{软}，则每个 $A$-\Term{module}{模}都\Term{soft-sheaf}{软}（T.F. II 3.7.1）。\Term{soft-sheaf}{软}\Term{sheaf-of-rings}{环层}的例子包括\Term{paracompact-space}{仿紧空间}上的\Term{continuous}{连续}实值或复值函数\Term{sheaf}{层}，以及 $C^r$ \Term{manifold}{流形}上的 $C^r$ 函数\Term{sheaf}{层}。
\begin{SGAresult}{命题}{2.3.2}
设 $S$ 为\Term{compact-space}{紧空间}，赋予一个\Term{soft-sheaf}{软}\Term{sheaf-of-rings}{环层} $A$。以 $\operatorname{Mod}(S)$ 表示左 $A$-\Term{module}{模}的\Term{category}{范畴}，置 $D(S)=D(\operatorname{Mod}(S))$。

a) $\operatorname{Lib}(S)$（相应地 $\operatorname{Loclib}(S)$，2.0）在 $\operatorname{Mod}(S)$ 中\Term{liftable}{可提升}，且对\Term{epimorphism}{满态射}的\Term{kernel}{核}\Term{quasi-stable}{拟稳定}（相应地\Term{stable}{稳定}，第 I 讲 1.2）。\footnote{更\Term{generic}{一般}地，$S$ \Term{paracompact-space}{仿紧}时此断言仍成立。}

b) 设 $E\in\operatorname{ob}D(S)$ 且 $n\in\mathbb Z$。$E$ 为 $n$-\Term{pseudo-coherent}{伪凝聚}（相应地\Term{pseudo-coherent}{伪凝聚}）的充要条件，是它相对于 $\operatorname{Lib}(S)$ 或 $\operatorname{Loclib}(S)$ 如此。\Term{canonical}{典范}\Term{functor}{函子}
\[
K^-(\operatorname{Lib}(S))\longrightarrow D(S),\qquad
K^-(\operatorname{Loclib}(S))\longrightarrow D(S)
\]
均\Term{fully-faithful}{全忠实}，\Term{essential-image}{本质像}为 $D^-(S)_{\mathrm{coh}}$（第 I 讲 2.15）。

c) 设 $E\in\operatorname{ob}D(S)$，$[m,n]\subset\mathbb Z$ 为区间。$E$ 的\Term{perfect-amplitude}{完美振幅}包含于 $[m,n]$（第 I 讲 4.8）中的充要条件，是存在\Term{quasi-isomorphism}{拟同构} $F\to E$，其中 $F$ 为\Term{complex}{复形}，满足 $F^i=0$ 对 $i\notin[m,n]$ 成立，$F^i\in\operatorname{ob}\operatorname{Lib}(S)$ 对 $i>m$ 成立，且 $F^m\in\operatorname{ob}\operatorname{Loclib}(S)$。\Term{canonical}{典范}\Term{functor}{函子}
\[
K^b(\operatorname{Loclib}(S))\longrightarrow D(S)
\]
是\Term{fully-faithful}{全忠实}的，\Term{essential-image}{本质像}为 $D(S)_{\mathrm{parf}}$（第 I 讲 4.9）。

d) 以 $\operatorname{Mod}(A(S))$ 表示左 $A(S)$-\Term{module}{模}的\Term{category}{范畴}，$D(A(S))$ 为相应的\Term{derived-category}{导出范畴}。\Term{functor}{函子}
\[
H^0(S,\,\cdot\,):\operatorname{Mod}(S)\longrightarrow\operatorname{Mod}(A(S))
\]
\Term{exact-functor}{正合}，且为\Term{equivalence-of-categories}{范畴等价}。此外，它定义\Term{equivalence-of-categories}{范畴等价}
\SourcePage{188}
\[
D(S)\longrightarrow D(A(S)),\qquad
D^-(S)_{\mathrm{coh}}\longrightarrow D^-(A(S))_{\mathrm{coh}},\qquad
D(S)_{\mathrm{parf}}\longrightarrow D(A(S))_{\mathrm{parf}},
\]
其中 $D(A(S))$ 中的\Term{pseudo-coherent}{伪凝聚}（相应地\Term{perfect-complex}{完美}）概念相对于\Term{finite-type}{有限型}\Term{free-module}{自由}（相应地\Term{projective-module}{投射}）\Term{module}{模}的\Term{category}{范畴}而言。它还给出 $D(S)$ 中由 $n$-\Term{pseudo-coherent}{伪凝聚}\Term{object}{对象}（相应地\Term{perfect-amplitude}{完美振幅}包含于 $[m,n]$ 的\Term{object}{对象}）组成的\Term{full-subcategory}{全子范畴}与 $D(A(S))$ 中相应\Term{full-subcategory}{全子范畴}之间的\Term{equivalence-of-categories}{等价}。
\end{SGAresult}
\begin{proof}
a) 的第一断言由 2.1.2 a) 得到：取 $\mathcal S$ 为 $S$ 上\Term{set}{集合}\Term{sheaf}{层}的\Term{topos}{拓扑斯}，$\mathcal C'_S=\operatorname{Mod}(S)$，并考虑 2.3.1。第二断言由 1.1 a)（相应地 1.2 a)）得到：取 $\mathcal C=\operatorname{Mod}$（变化的\Term{open-subset}{开集}上 $A$-\Term{module}{模}的\Term{fibered-category}{纤维范畴}），$\mathcal C_0=\operatorname{Lib}$（相应地 $\operatorname{Loclib}$）；由 2.3.1 和 2.1.1 b)，1.1 的条件 (ii) 成立。b) 应用 1.1 b)、c)（$\mathcal C,\mathcal C_0$ 如上）。c) 由 1.2 b)、c)、d) 得到，其中 $\mathcal C=\operatorname{Mod}$，$\mathcal C_0=\operatorname{Lib}$。d) 的第一断言由 2.3.1 和 2.1.2 b)(ii) 得到，取 $\mathcal C'_S=\mathcal C_S=\operatorname{Mod}(S)$。d) 的其余断言由第一断言、b)、c) 以及如下事实得到：$\Gamma_S$ 将\Term{finite-type}{有限型}\Term{free-module}{自由}\Term{module}{模}（相应地\Term{free-module}{自由模}的\Term{direct-summand}{直和因子}）变为\Term{finite-type}{有限型}\Term{free-module}{自由}（相应地\Term{projective-module}{投射}）\Term{module}{模}。
\end{proof}
\textbf{2.4。Stein \Term{space}{空间}。}
\begin{SGAresult}{定义}{2.4.1}
设 $\mathcal S$ 为\Term{locally-ringed-topos}{局部赋环拓扑斯}（第 I 讲 2.15.1 (ii)），\Term{final-object}{终对象}为 $S$。称 $\mathcal S$ 为 \Term{stein-topos}{Stein 拓扑斯}，若\Term{structure-sheaf}{结构层} $\mathcal O_S$ \Term{coherent}{凝聚}（第 I 讲 3.1），且有：

“\Term{theorem}{定理} A”：每个\Term{coherent}{凝聚} $\mathcal O_S$-\Term{module}{模}（第 I 讲 3.1）由其\Term{section}{截面}生成。

“\Term{theorem}{定理} B”：对每个\Term{coherent}{凝聚} $\mathcal O_S$-\Term{module}{模} $F$，$H^i(S,F)=0$ 对 $i>0$ 成立。
\end{SGAresult}
\SourcePage{189}
\textbf{注记 2.4.1.1。}若 $\mathcal S$ 有“足够多的\Term{closed-point}{闭点}” $s$，\footnote{即有足够多形如 $i_s:s\to\mathcal S$ 的\Term{fiber-functor}{纤维函子}（SGA 4 IV 6.4.1），其中 $i_s$ 为 $\mathcal S$ 的一个闭\Term{subobject}{子对象}的包含。}使 $s$ 所定义的\Term{ideal}{理想} $\mathfrak m_s$ \Term{coherent}{凝聚}，\footnote{即由 $\mathcal O_S$ 中那些在 $s$ 处的\Term{stalk}{茎}属于 $\mathcal O_{S,s}$ 的\Term{maximal-ideal}{极大理想} $m_s$ 的\Term{section}{截面}组成。}则条件“对每个\Term{coherent}{凝聚}\Term{sheaf}{层} $F$，$H^1(S,F)=0$”蕴含\Term{theorem}{定理} A。事实上，设 $F$ 为\Term{coherent}{凝聚} $\mathcal O_S$-\Term{module}{模}，$s$ 为 $\mathcal S$ 的\Term{closed-point}{闭点}。有\Term{exact-sequence}{正合列}
\[
0\longrightarrow\mathfrak m_sF\longrightarrow F\longrightarrow F/\mathfrak m_sF\longrightarrow0,
\]
其中 $F/\mathfrak m_sF$ 为集中在 $s$ 上的\Term{coherent}{凝聚}\Term{sheaf}{层}。应用\Term{functor}{函子} $H^0(S,\,\cdot\,)$，写出\Term{cohomology}{上同调}\Term{exact-sequence}{正合列}，得 $H^0(S,F)\to H^0(S,F/\mathfrak m_sF)$ 为\Term{epimorphism}{满态射}。由于 $H^0(S,F/\mathfrak m_sF)\simeq F_s/m_sF_s$，其中 $m_s$ 为 $\mathcal O_{S,s}$ 的\Term{maximal-ideal}{极大理想}，由 \Term{nakayama-lemma}{Nakayama 引理}即得结论。

撰稿人不知道上述蕴含在\Term{generic}{一般}情形下是否成立（只假设 $\mathcal O_S$ 为\Term{coherent}{凝聚}的\Term{local-ring}{局部环}\Term{sheaf}{层}）。

撰稿人也不知道条件“对每个\Term{coherent}{凝聚}\Term{sheaf}{层} $F$，$H^1(S,F)=0$”是否蕴含\Term{theorem}{定理} B。

\textbf{例 2.4.2。}复 Stein \Term{analytic-space}{解析空间} [3]、赋予\Term{holomorphic}{全纯}函数\Term{sheaf}{层}所诱导的\Term{sheaf}{层}的 $\mathbb C^n$ 中的凸紧集 [3]、\Term{noetherian-ring}{诺特环}的\Term{spectrum}{谱}、\Term{affine-scheme}{仿射}\Term{rigidity}{刚性}\Term{analytic-space}{解析空间} [13]，都定义 \Term{stein-topos}{Stein 拓扑斯}。
\begin{SGAresult}{\Term{lemma}{引理}}{2.4.3}
设 $\mathcal S$ 为 \Term{stein-topos}{Stein 拓扑斯}（2.4.1）。假定\Term{final-object}{终对象} $S$ \Term{quasi-compact}{拟紧}（2.1.1），以 $\operatorname{Coh}(S)$ 表示\Term{coherent}{凝聚}，即\Term{finite-presentation}{有限表现} $\mathcal O_S$-\Term{module}{模}的\Term{abelian-category}{阿贝尔范畴}。则：

a) $\operatorname{Lib}(S)$（相应地 $\operatorname{Loclib}(S)$，2.0）在 $\operatorname{Coh}(S)$ 中\Term{liftable}{可提升}，且对\Term{epimorphism}{满态射}的\Term{kernel}{核}\Term{quasi-stable}{拟稳定}（相应地\Term{stable}{稳定}，第 I 讲 1.2）。

b) \Term{canonical}{典范}\Term{functor}{函子}
\[
K^-(\operatorname{Lib}(S))\longrightarrow K^-(\operatorname{Loclib}(S))\longrightarrow D^-(\operatorname{Coh}(S))
\]
均为\Term{equivalence-of-categories}{范畴等价}。
\SourcePage{190}
c) 置 $A=\Gamma(S,\mathcal O_S)$，以 $\operatorname{Mod}(A)$ 表示 $A$-\Term{module}{模}的\Term{category}{范畴}，$D(A)$ 为相应的\Term{derived-category}{导出范畴}。\Term{functor}{函子}
\[
\Gamma_S:\operatorname{Coh}(S)\longrightarrow\operatorname{Mod}(A)
\]
\Term{exact-functor}{正合}且\Term{fully-faithful}{全忠实}，\Term{essential-image}{本质像}为 $\operatorname{Mod}(A)$ 中由\Term{finite-presentation}{有限表现} $A$-\Term{module}{模}组成的\Term{strictly-full-subcategory}{严格全}\Term{additive-category}{加性子范畴}。它诱导\Term{fully-faithful}{全忠实}\Term{functor}{函子}
\[
D^-(\operatorname{Coh}(S))\longrightarrow D(A),
\]
其\Term{essential-image}{本质像}为由相对于\Term{finite-type}{有限型}\Term{free-module}{自由}\Term{module}{模}的\Term{subcategory}{子范畴}\Term{pseudo-coherent}{伪凝聚}的\Term{complex}{复形}组成的\Term{full-subcategory}{全子范畴} $D^-(A)_{\mathrm{coh}}$。
\end{SGAresult}
\begin{proof}
对 a) 的第一断言，取 $\mathcal C'_S=\operatorname{Coh}(S)$ 应用 2.1.2 a)；\Term{theorem}{定理} B 保证 2.1.2 的假设成立。沿用 2.1.2 a) 的记号，若 $E$ \Term{coherent}{凝聚}而 $G$ 为\Term{finite-type}{有限型}\Term{locally-free}{局部自由}，则 $\mathcal Hom(G,E)$ \Term{coherent}{凝聚}；故 $\operatorname{Loclib}(S)$（从而 $\operatorname{Lib}(S)$）在 $\operatorname{Coh}(S)$ 中\Term{liftable}{可提升}。a) 的第二断言由 1.1 a)（相应地 1.2 a)）得到，取 $\mathcal C_0=\operatorname{Lib}$（相应地 $\operatorname{Loclib}$）；\footnote{由第 I 讲 2.15.1，$\operatorname{Lib}$ 和 $\operatorname{Loclib}$ 对\Term{epimorphism}{满态射}的\Term{kernel}{核}\Term{local-scheme}{局部}\Term{stable}{稳定}。}由\Term{theorem}{定理} A 和 2.1.1 b)，1.1 的条件 (ii) 中\Term{lie-bracket}{括号}内的版本成立。依次取 $\mathcal C_0=\operatorname{Lib}$、$\mathcal C_0=\operatorname{Loclib}$，由 1.1 c) 得到 b)。c) 的第一断言由 2.1.2 b)(ii) 对 $\mathcal C'_S=\operatorname{Coh}(S)$ 应用得到。为证明 $\Gamma_S:D^-(\operatorname{Coh}(S))\to D(A)$ \Term{fully-faithful}{全忠实}，由 b)，只须证明复合
\[
K^-(\operatorname{Lib}(S))\longrightarrow D^-(\operatorname{Coh}(S))\xrightarrow{\Gamma_S}D(A)
\]
\Term{fully-faithful}{全忠实}。它\Term{resolution}{分解}为
\[
K^-(\operatorname{Lib}(S))\xrightarrow{(*)}K^-(\operatorname{Lib}(A))\xrightarrow{(**)}D(A),
\]
\SourcePage{191}
其中 $\operatorname{Lib}(A)$ 为\Term{finite-type}{有限型}\Term{free-module}{自由} $A$-\Term{module}{模}的\Term{category}{范畴}，(*) 由 $\Gamma_S$ 诱导。显然 (*) 为\Term{equivalence-of-categories}{范畴等价}。另一方面，取 $\mathcal C=\operatorname{Mod}(A)$、$\mathcal C_0=\operatorname{Lib}(A)$ 应用第 I 讲 2.7 b)，并考虑 $K^{-,\varnothing}(\operatorname{Lib}(A))=0$，得到 (**) \Term{fully-faithful}{全忠实}，其\Term{essential-image}{本质像}为 $D^-(A)_{\mathrm{coh}}$。这证明了 c) 的最后断言，也完成了 2.4.3 的证明。
\end{proof}
\begin{SGAresult}{\Term{lemma}{引理}}{2.4.4}
在 2.4.3 的假设下，\Term{canonical}{典范}\Term{functor}{函子}
\[
D^b(\operatorname{Coh}(S))\longrightarrow D(S)
\]
是\Term{fully-faithful}{全忠实}的，其\Term{essential-image}{本质像}为 $D^b(S)_{\mathrm{coh}}$（第 I 讲 2.15）。
\end{SGAresult}
\begin{proof}
取 $\mathcal A=\operatorname{Coh}(S)$、$\mathcal B=\operatorname{Mod}(S)$，$\phi:\mathcal A\to\mathcal B$ 为包含\Term{functor}{函子}，应用 2.1.3。因 $\phi$ \Term{fully-faithful}{全忠实}，由 2.1.3 a)(ii)，只须证明：对固定 $M\in\operatorname{ob}\operatorname{Coh}(S)$ 和 $n>0$，\Term{functor}{函子} $L\mapsto\operatorname{Ext}_{\operatorname{Mod}(S)}^n(L,M)$ 对 $L\in\operatorname{ob}\operatorname{Coh}(S)$ 是可余消去的。由\Term{theorem}{定理} A 及 2.1.1 b)，每个\Term{coherent}{凝聚} $\mathcal O_S$-\Term{module}{模}都是\Term{finite-type}{有限型}\Term{free-module}{自由模}的\Term{quotient-module}{商模}；而若 $L$ 为\Term{finite-type}{有限型}\Term{free-module}{自由}，则由\Term{theorem}{定理} B，$\operatorname{Ext}_{\operatorname{Mod}(S)}^n(L,M)=0$，断言得证。因此 $D^b(\operatorname{Coh}(S))\to D(S)$ \Term{fully-faithful}{全忠实}，\Term{essential-image}{本质像}为由\Term{cohomology}{上同调}\Term{bounded}{有界}且\Term{coherent}{凝聚}的\Term{complex}{复形}组成的\Term{full-subcategory}{全子范畴}；因 $\mathcal O_S$ \Term{coherent}{凝聚}，这正是 $D^b(S)_{\mathrm{coh}}$（第 I 讲 3.5 b)）。
\end{proof}
\begin{SGAresult}{命题}{2.4.5}
沿用 2.4.3 的假设和记号，有：

a) \Term{canonical}{典范}\Term{functor}{函子}
\[
K^{-,b}(\operatorname{Lib}(S))\longrightarrow K^{-,b}(\operatorname{Loclib}(S))\longrightarrow D^b(S)_{\mathrm{coh}}
\]
均为\Term{equivalence-of-categories}{范畴等价}。此外，$\mathbf R\Gamma_S:D^+(S)\to D^+(A)$ 诱导\Term{fully-faithful}{全忠实}\Term{functor}{函子}
\[
\mathbf R\Gamma_S:D^b(S)_{\mathrm{coh}}\longrightarrow D^+(A),
\]
\SourcePage{192}
其\Term{essential-image}{本质像}为 $D^b(A)_{\mathrm{coh}}$（$D(A)$ 中由\Term{cohomology}{上同调}\Term{bounded}{有界}且相对于\Term{finite-type}{有限型}\Term{free-module}{自由}\Term{module}{模}的\Term{category}{范畴}\Term{pseudo-coherent}{伪凝聚}的\Term{complex}{复形}组成的\Term{triangulated-category}{三角子范畴}）。

b) 设 $E\in\operatorname{ob}D(S)$，$[m,n]\subset\mathbb Z$ 为区间。若 $E$ 的\Term{perfect-amplitude}{完美振幅}包含于 $[m,n]$（第 I 讲 4.8），则 $E$ \Term{isomorphism}{同构}于一个\Term{complex}{复形} $F$，使 $F^i=0$ 对 $i\notin[m,n]$ 成立，$F^i\in\operatorname{ob}\operatorname{Lib}(S)$ 对 $i>m$ 成立，且 $F^m\in\operatorname{ob}\operatorname{Loclib}(S)$。

c) \Term{canonical}{典范}\Term{functor}{函子}
\[
K^b(\operatorname{Loclib}(S))\longrightarrow D(S)_{\mathrm{parf}}\quad(\text{第 I 讲 }4.9)
\]
为\Term{equivalence-of-categories}{范畴等价}。此外，$\mathbf R\Gamma_S:D^+(S)\to D^+(A)$ 诱导\Term{fully-faithful}{全忠实}\Term{functor}{函子}
\[
\mathbf R\Gamma_S:D(S)_{\mathrm{parf}}\longrightarrow D^+(A),
\]
其\Term{essential-image}{本质像}为 $D(A)_{\mathrm{parf}}$（$D(A)$ 中由相对于\Term{finite-type}{有限型}\Term{projective-module}{投射}\Term{module}{模}的\Term{category}{范畴}\Term{perfect-complex}{完美}的\Term{complex}{复形}组成的\Term{triangulated-category}{三角子范畴}）。
\end{SGAresult}
\begin{proof}
a) 由 2.4.3 b)、c)，\Term{functor}{函子}
\[
K^{-,b}(\operatorname{Lib}(S))\longrightarrow K^{-,b}(\operatorname{Loclib}(S))\longrightarrow D^b(\operatorname{Coh}(S))\xrightarrow{\Gamma_S}D^b(A)_{\mathrm{coh}}
\]
均为\Term{equivalence-of-categories}{范畴等价}。因\Term{fully-faithful}{全忠实}\Term{functor}{函子} $\Gamma_S:D^b(\operatorname{Coh}(S))\to D(A)$ \Term{resolution}{分解}为
\[
D^b(\operatorname{Coh}(S))\longrightarrow D^b(S)_{\mathrm{coh}}\xrightarrow{\mathbf R\Gamma_S}D(A),
\]
由 2.4.4 即得结论。

b) 假定 $\operatorname{parf.amp}(E)\subset[m,n]$。特别地，$E$ \Term{pseudo-coherent}{伪凝聚}；由 a)，以\Term{isomorphism}{同构}\Term{object}{对象}替代后，可假定 $E^i\in\operatorname{ob}\operatorname{Lib}(S)$ 对 $i>m$ 成立，且 $E^i=0$ 对 $i>n$ 成立。
\SourcePage{193}
作\Term{truncation}{截断}后，甚至可假定 $E^i=0$ 对 $i<m$ 成立。取 $\mathcal C=\operatorname{Mod}$、$\mathcal C_0=\operatorname{Lib}$ 应用第 I 讲 4.16，即得 $E^m\in\operatorname{ob}\operatorname{Loclib}(S)$，故断言成立。

c) 由 2.4.3 a)（相应地 2.1.1 b) 和\Term{theorem}{定理} A），对 $\mathcal C_S=\operatorname{Coh}(S)$、$\mathcal C_0=\operatorname{Lib}$，1.1 (i)（相应地 1.1 (ii) 中\Term{lie-bracket}{括号}内的版本）成立。因此由 1.2 c)、d)，\Term{functor}{函子}
\[
K^b(\operatorname{Loclib}(S))\longrightarrow D^b(\operatorname{Coh}(S))
\]
\Term{fully-faithful}{全忠实}。由 2.4.4，\Term{functor}{函子}
\[
K^b(\operatorname{Loclib}(S))\longrightarrow D(S)
\]
也\Term{fully-faithful}{全忠实}，由 b)，其\Term{essential-image}{本质像}为 $D(S)_{\mathrm{parf}}$。仿照 2.4.3 c) 的证明可知，由 $\Gamma_S$ 诱导的\Term{functor}{函子} $K^b(\operatorname{Loclib}(S))\to D(S)$ \Term{fully-faithful}{全忠实}，\Term{essential-image}{本质像}为 $D(S)_{\mathrm{parf}}$；最后如 a) 一样完成证明。
\end{proof}
\textbf{注记 2.4.6。}撰稿人不知道 2.4.4 中是否可将 $D^b$ 替换为 $D^-$，从而也不知道在将 $K^{-,b}$ 替换为 $K^-$、$D^b$ 替换为 $D^-$ 后，2.4.5 a) 是否仍成立。

最后给出 2.4.5 的一个有趣补充。先给出定义：
\begin{SGAresult}{定义}{2.4.7}
称一个\Term{ringed-space}{赋环空间}为\Term{complex-pseudoanalytic-space}{复伪解析空间}（espace pseudo-analytique complexe），若它\Term{local-scheme}{局部}\Term{isomorphism}{同构}于形如 $(Y,\mathcal O_X|Y)$ 的\Term{ringed-space}{赋环空间}，其中 $X$ 为复\Term{analytic-space}{解析空间}，$Y$ 为 $X$ 的\Term{closed-subset}{闭子集}。

因此，\Term{complex-pseudoanalytic-space}{复伪解析空间}的\Term{structure-sheaf}{结构层}是\Term{coherent}{凝聚}的。

称\Term{complex-pseudoanalytic-space}{复伪解析空间}为 Stein 的，若其对应\Term{topos}{拓扑斯}为 \Term{stein-topos}{Stein 拓扑斯}（2.4.1）。
\end{SGAresult}
\SourcePage{194}
\begin{SGAresult}{命题}{2.4.8}
\footnote{A. Grothendieck 将这一结果及其证明慷慨地告知撰稿人；J. Frisch [5] 也独立得到了它。}设 $S$ 为紧的、Stein 的\Term{complex-pseudoanalytic-space}{复伪解析空间}（2.4.7）。假定对 $S$ 的每个\Term{open-subset}{开集} $U$ 及 $U$ 的每个\Term{reduced}{约化}\Term{closed-subspace}{闭子空间} $X$，$X$ 的\Term{normalization}{正规化}是\Term{locally-connected}{局部连通}的。\footnote{这里承认 [10] 中\Term{local-scheme}{局部}\Term{analytic-geometry}{解析几何}的概念可推广到\Term{complex-pseudoanalytic-space}{复伪解析空间}；事实上，每一次都可直接归结到通常情形。}（例如，若 $S$ \Term{local-scheme}{局部}由 $\mathbb R^n$ 中的实解析等式或不等式定义，则此条件成立 [15]。）则 $A=\Gamma(S,\mathcal O_S)$ 为\Term{noetherian-ring}{诺特环}。
\end{SGAresult}
\begin{proof}
须证：若 $M$ 为\Term{finite-presentation}{有限表现} $A$-\Term{module}{模}，则 $M$ 的\Term{submodule}{子模}全体（\Term{equivalence}{等价}地，\Term{finite-type}{有限型}\Term{submodule}{子模}全体）满足极大条件。若 $N$ 为 $M$ 的\Term{finite-type}{有限型}\Term{submodule}{子模}，则由第 I 讲 2.6、2.9，$M/N$ 为\Term{finite-presentation}{有限表现}，由 2.4.3 c) 的\Term{equivalence-of-categories}{等价}，$N$ 也为\Term{finite-presentation}{有限表现}。再次使用 2.4.3 c) 中\Term{coherent}{凝聚} $\mathcal O_S$-\Term{module}{模}与\Term{finite-presentation}{有限表现} $A$-\Term{module}{模}之间的\Term{equivalence-of-categories}{等价}及 $S$ 的紧性，即归结到证明下述\Term{lemma}{引理}。
\end{proof}
\begin{SGAresult}{\Term{lemma}{引理}}{2.4.8.1}
设 $S$ 为满足 2.4.8 中\Term{locally-connected}{局部连通}性条件的\Term{complex-pseudoanalytic-space}{复伪解析空间}。设 $F$ 为\Term{coherent}{凝聚} $\mathcal O_S$-\Term{module}{模}，$(E_n)$ 为 $F$ 的\Term{coherent}{凝聚} $\mathcal O_S$-\Term{submodule}{子模}的递增序列，$s\in S$。则存在 $s$ 的\Term{open-neighborhood}{开邻域} $U$，使序列 $(E_n|U)$ 最终\Term{invariant}{不变}。
\end{SGAresult}
\begin{proof}
环 $\mathcal O_{S,s}$ \Term{noetherian-ring}{诺特}，故存在整数 $n_0$，使 $(E_n)_s=(E_{n_0})_s$ 对 $n\geq n_0$ 成立。将 $(E_n)$ 替换为 $(E_n/E_{n_0})_{n\geq n_0}$，可假定 $(E_n)_s=0$ 对每个 $n$ 成立。另一方面，$F_s=M$ 为\Term{finite-type}{有限型} $\mathcal O_s$-\Term{module}{模}，故存在\Term{finite}{有限}\Term{filtration}{滤过}（Bourbaki，Alg. com.，第 IV 章 §1，\Term{theorem}{定理} 1）
\[
M=M_0\supset M_1\supset\cdots\supset M_q=0,
\]
使 $M_i/M_{i+1}\simeq\mathcal O_s/\mathfrak p_i$，其中 $\mathfrak p_i$ 为 $\mathcal O_s$ 的\Term{prime-ideal}{素理想}。因\Term{finite-type}{有限型} $\mathcal O_s$-\Term{module}{模}的\Term{category}{范畴}与\Term{coherent}{凝聚} $\mathcal O_S$-\Term{module}{模}在 $s$ 处的\Term{germ}{芽}的\Term{category}{范畴}\Term{equivalence-of-categories}{等价}，缩小到 $s$ 的\Term{neighborhood}{邻域}后，可假定
\SourcePage{195}
$M_i$（相应地 $\mathfrak p_i$）是 $F$（相应地 $\mathcal O_S$）的\Term{coherent}{凝聚}\Term{submodule}{子模} $F_i$（相应地\Term{coherent}{凝聚}\Term{ideal}{理想} $\mathcal P_i$）在 $s$ 处的\Term{stalk}{茎}，满足
\[
F=F_0\supset F_1\supset\cdots\supset F_q=0,\qquad F_i/F_{i+1}\simeq\mathcal O_S/\mathcal P_i.
\]
证明 $(E_n)$ 在 $s$ \Term{neighborhood}{邻域}上最终\Term{invariant}{不变}，显然\Term{equivalence}{等价}于证明它在 $F$ 的\Term{filtration}{滤过} $(F_i)$ 的各个\Term{graded}{分次}商上诱导的序列最终\Term{invariant}{不变}。因此可假定 $F=\mathcal O_S$，且 $\mathcal O_{S,s}$ 为\Term{integral-domain}{整环}，各 $E_n$ 是 $\mathcal O_S$ 的递增\Term{ideal}{理想}序列，满足 $(E_n)_s=0$。因 $\mathcal O_s$ \Term{integral-domain}{整}，特别地\Term{reduced}{约化}，由 Oka \Term{theorem}{定理}（[10] IV §3），缩小到 $s$ \Term{neighborhood}{邻域}后，还可假定 $S$ \Term{reduced}{约化}。

证明存在 $s$ 的\Term{open-neighborhood}{开邻域} $U$，使充分大的 $n$ 满足 $E_n|U=0$。考虑 $S$ 的\Term{normalization}{正规化} $\widetilde S$（[10] IV §4）。已知（同上）$\widetilde S$ 是\Term{finite}{有限} $\mathcal O_S$-\Term{algebra}{代数} $\widetilde{\mathcal O}_S$ 的\Term{analytic-spectrum}{解析谱}（[10] II §3），其在每个 $t\in S$ 处的\Term{stalk}{茎}为 $\mathcal O_{S,t}$ 的\Term{integral-closure}{整闭包}（参见 2.4.8 脚注 2）。有 $E_n\subset E_n\widetilde{\mathcal O}_S$，故只须证明充分大的 $n$ 满足 $E_n\widetilde{\mathcal O}_S=0$。又有（[10] II §3，命题 8）
\[
E_n\widetilde{\mathcal O}_S=p_*\bigl(E_n\mathcal O_{\widetilde S}\bigr),
\]
其中 $p:\widetilde S\to S$ 为\Term{canonical-projection}{典范投影}，$E_n\mathcal O_{\widetilde S}$ 为 $E_n\widetilde{\mathcal O}_S$ \Term{canonical}{典范}定义的 $\mathcal O_{\widetilde S}$-\Term{ideal}{理想}。因此，只须证明存在 $s$ 在 $S$ 中的\Term{open-neighborhood}{开邻域} $V$，使充分大的 $n$ 满足 $E_n\widetilde{\mathcal O}_S|U=0$；这将问题归结到 $S$ \Term{normal}{正规}的情形，以下便作此假设。

因 $(E_n)_s=0$，由 $E_n$ 定义的\Term{closed-immersion}{闭浸入} $i_n:S_n\to S$ 在 $s$ 处\Term{etale}{平展}。又因 $S$ \Term{normal}{正规}，$i_n$ \Term{etale}{平展}的点集既开又闭（见下面的\Term{lemma}{引理}），故其在 $S$ 中的像也是既开又闭的\Term{subset}{子集}。因此 $E_n$ 在 $s$ 所在的\Term{connected-component}{连通分支}上的\Term{restriction}{限制}为零；由 $S$ \Term{locally-connected}{局部连通}，此\Term{connected-component}{连通分支}为\Term{open-subset}{开集}，结论遂得。最后只须证明：
\end{proof}
\SourcePage{196}
\begin{SGAresult}{\Term{lemma}{引理}}{2.4.3.2}
（参见 EGA IV 18.10.3。）设 $Y$ 为\Term{normal}{正规}\Term{pseudoanalytic-space}{伪解析空间}，\footnote{只须 $Y$ 的\Term{local-ring}{局部环}均为\Term{integral-domain}{整环}。}$f:X\to Y$ 为\Term{unramified}{非分歧}\Term{morphism}{态射}。$X$ 中 $f$ \Term{etale}{平展}的点集既开又闭。
\end{SGAresult}
\begin{proof}
$f$ \Term{etale}{平展}的点集显然是\Term{open-subset}{开集}（因为 $f$ 在某点\Term{etale}{平展}，意味着它在该点\Term{neighborhood}{邻域}上是\Term{isomorphism}{同构}，[8] VI）。还须证明 $f$ 不\Term{etale}{平展}的点集为\Term{open-subset}{开集}。因\Term{unramified}{非分歧}\Term{morphism}{态射}\Term{local-scheme}{局部}为\Term{immersion}{浸入}（[8] VI），可假定 $f$ 为由\Term{coherent}{凝聚}\Term{ideal}{理想} $I$ 定义的\Term{closed-immersion}{闭浸入}。设 $x\in X$，且 $I_x\ne0$。取 $\mathcal O_{Y,x}$ 的一个包含 $I_x$ 的\Term{prime-ideal}{素理想} $\mathfrak p$。有 $\dim(\mathcal O_{Y,x})_{\mathfrak p}>0$，否则，由 $\mathcal O_{Y,x}$ 为\Term{integral-domain}{整环}，将有 $\mathfrak p=0$。故 $Y$ 在 $x$ 处的\Term{codimension}{余维} $>0$，进而在 $x$ 的某\Term{neighborhood}{邻域}中\Term{codimension}{余维} $>0$（[10] III A，\Term{lemma}{引理} 3 的推论 1），结论即得。
\end{proof}
\SGAsection{3}{关于概形上拟凝聚层的补充}
\footnote{本节主要归功于 J. L. Verdier。}
\textbf{3.0。}本节固定一个\Term{quasi-compact}{拟紧}\Term{quasi-separated}{拟分离}\Term{scheme}{概形} $S$。\Term{canonical}{典范}的（\Term{fully-faithful}{全忠实}）包含\Term{functor}{函子}
\[
\phi:\operatorname{Qcoh}(S)\longrightarrow\operatorname{Mod}(S)
\]
\footnote{$\operatorname{Mod}(S)$（相应地 $\operatorname{Qcoh}(S)$）表示 $\mathcal O_S$-\Term{module}{模}（相应地\Term{quasi-coherent}{拟凝聚} $\mathcal O_S$-\Term{module}{模}）的\Term{category}{范畴}。}延拓为从 $D(\operatorname{Qcoh}(S))$ 到 $D(S)=D(\operatorname{Mod}(S))$ 的\Term{functor}{函子}，仍记为 $\phi$；本节研究这个\Term{functor}{函子}。\footnote{在不致引起混淆的公式中，将省略 $\phi$。}
\begin{SGAresult}{\Term{lemma}{引理}}{3.1}
$\operatorname{Qcoh}(S)$ 是满足 (AB 5) 的\Term{abelian-category}{阿贝尔范畴}，其中\Term{finite-type}{有限型}\Term{quasi-coherent}{拟凝聚}\Term{sheaf}{层}构成一个“小”的\Term{generators}{生成元族}。特别地（Tohoku 1.10.1），$\operatorname{Qcoh}(S)$ 有足够多\Term{injective-object}{内射对象}。
\end{SGAresult}
\SourcePage{197}
\begin{proof}
关于\Term{generators}{生成元}的断言，立即由下述事实得到：$S$ 上每个\Term{quasi-coherent}{拟凝聚}\Term{sheaf}{层}都是其\Term{finite-type}{有限型}\Term{quasi-coherent}{拟凝聚}\Term{subsheaf}{子层}的\Term{filtered-inductive-limit}{滤过归纳极限}（EGA I 9.4.9、IV 1.7.7）。其余断言显然。
\end{proof}
\textbf{注记 3.1.1。}若只假设 $S$ \Term{quasi-separated}{拟分离}，撰稿人不知道 $\operatorname{Qcoh}(S)$ 是否仍有足够多\Term{injective-object}{内射对象}。
\begin{SGAresult}{\Term{lemma}{引理}}{3.2}
\Term{functor}{函子} $\phi:\operatorname{Qcoh}(S)\to\operatorname{Mod}(S)$ 有\Term{right-adjoint}{右伴随}
\[
Q:\operatorname{Mod}(S)\longrightarrow\operatorname{Qcoh}(S),
\]
称为\Term{coherator}{凝聚化函子}（cohéreur）。
\end{SGAresult}
\begin{proof}
若 $S$ \Term{affine-scheme}{仿射}，则 $Q$ 的存在是直接的：取 $Q=\widetilde{\Gamma_S}$（EGA $\operatorname{Err}_{\mathrm I}$ 1.3.7）。现设 $f:U\to S$ 为 $S$ 的\Term{affine-scheme}{仿射}\Term{open-subset}{开集}。显然，对 $E\in\operatorname{ob}\operatorname{Qcoh}(S)$ 及 $F\in\operatorname{ob}\operatorname{Mod}(U)$，有\Term{canonical}{典范}\Term{isomorphism}{同构}
\[
\operatorname{Hom}(E,f_*F)\simeq\operatorname{Hom}(E,f_*Q(F)).\tag{1}
\]
取 $S$ 的\Term{finite}{有限}\Term{affine-scheme}{仿射}开\Term{covering}{覆盖} $(f_i:U_i\to S)$；对每个 $(i,j)$，取 $U_{ij}=U_i\cap U_j$ 的\Term{finite}{有限}\Term{affine-scheme}{仿射}开\Term{covering}{覆盖} $(U_{ijk}\to U_{ij})$（因 $S$ \Term{quasi-separated}{拟分离}，$U_{ij}$ \Term{quasi-compact}{拟紧}）。设 $F\in\operatorname{ob}\operatorname{Mod}(S)$。有\Term{exact-sequence}{正合列}
\[
0\longrightarrow F\longrightarrow\prod_i f_{i*}(F|U_i)\longrightarrow\prod_{ijk}f_{ijk*}(F|U_{ijk}),\tag{2}
\]
其中 $f_{ijk}$ 表示 $U_{ijk}$ 到 $S$ 的包含。注意，若 $g:V\to U$ 为 $S$ 的\Term{affine-scheme}{仿射}\Term{open-subset}{开集}之间的包含，则有显然的自然\Term{arrow}{箭头} $Q(F|U)\to g_*Q(F|V)$。于是从 (2) 得到\Term{arrow}{箭头}
\[
\prod_i f_{i*}Q(F|U_i)\longrightarrow\prod_{ijk}f_{ijk*}Q(F|U_{ijk}),\tag{3}
\]
以 $QF$ 表示其\Term{kernel}{核}。设 $E\in\operatorname{ob}\operatorname{Qcoh}(S)$，对 (2) 应用\Term{functor}{函子} $\operatorname{Hom}(E,\,\cdot\,)$，结合 (1)，得关于 $E$ \Term{functorial}{函子性}的\Term{isomorphism}{同构}
\SourcePage{198}
\[
\operatorname{Hom}(E,F)\simeq\operatorname{Hom}(E,QF).
\]
这证明了 $Q$ 的存在；因它为\Term{adjoint-functor}{伴随函子}，唯一性显然。
\end{proof}
\begin{SGAresult}{\Term{lemma}{引理}}{3.3}
a) \Term{coherator}{凝聚化函子}是将\Term{injective-object}{内射对象}变为\Term{injective-object}{内射对象}的\Term{additive-functor}{加性函子}。

b) 以
\[
\mathbf RQ:D^+(S)\longrightarrow D^+(\operatorname{Qcoh}(S))
\]
表示\Term{coherator}{凝聚化函子}的\Term{right-derived-functor}{右导出函子}。对 $E\in\operatorname{ob}D(\operatorname{Qcoh}(S))$、$F\in\operatorname{ob}D^+(S)$，有\Term{canonical}{典范}\Term{functorial}{函子性}\Term{isomorphism}{同构}
\[
\mathbf R\operatorname{Hom}_{\mathrm{Mod}}(E,F)\simeq\mathbf R\operatorname{Hom}_{\mathrm{Qcoh}}(E,\mathbf RQ(F)),\tag{i}
\]
\[
\operatorname{Hom}_{\mathrm{Mod}}(E,F)\simeq\operatorname{Hom}_{\mathrm{Qcoh}}(E,\mathbf RQ(F)).\tag{ii}
\]
\footnote{下标 $\mathrm{Mod}$（相应地 $\mathrm{Qcoh}$）是 $D(S)$（相应地 $D(\operatorname{Qcoh}(S))$）的简写。}
\end{SGAresult}
\begin{proof}
a) 立即由定义得到。为证 b)，只须建立 (i)，因为取\Term{cohomology}{上同调}便得到 (ii)。可假定 $F$ 的各\Term{set-component}{分量}\Term{injective-object}{内射}；由 a)，$QF$ 的各\Term{set-component}{分量}也\Term{injective-object}{内射}，于是 (i) 直接来自 $Q$ 的定义。
\end{proof}
\begin{SGAresult}{\Term{lemma}{引理}}{3.4}
设 $f:X\to Y$ 为\Term{quasi-compact}{拟紧}\Term{quasi-separated}{拟分离}\Term{scheme}{概形}间的\Term{morphism}{态射}。对 $E\in\operatorname{ob}\operatorname{Mod}(X)$，有\Term{canonical}{典范}\Term{functorial}{函子性}\Term{isomorphism}{同构}
\[
Qf_*(E)\simeq f_*Q(E).
\]
\end{SGAresult}
\SourcePage{199}
\begin{proof}
因 $f$ \Term{quasi-compact}{拟紧}\Term{quasi-separated}{拟分离}（EGA IV 1.2.3、1.2.4），$f_*$ 将\Term{quasi-coherent}{拟凝聚}\Term{sheaf}{层}变为\Term{quasi-coherent}{拟凝聚}\Term{sheaf}{层}（III 1.4.10、IV 1.7.21）。从 $\operatorname{Mod}(X)$ 到 $\operatorname{Qcoh}(Y)$ 的\Term{functor}{函子} $Qf_*$ 和 $f_*Q$ 都是 $f^*:\operatorname{Qcoh}(Y)\to\operatorname{Mod}(X)$ 的\Term{right-adjoint}{右伴随}，故\Term{canonical}{典范}\Term{isomorphism}{同构}。
\end{proof}
\begin{SGAresult}{命题}{3.5}
若 $S$ \Term{noetherian-scheme}{诺特}，或 $S$ \Term{quasi-compact}{拟紧}且\Term{separated}{分离}，则\Term{functor}{函子}
\[
\phi:D^+(\operatorname{Qcoh}(S))\longrightarrow D^+(S)\qquad(3.0)
\]
\Term{fully-faithful}{全忠实}，其\Term{essential-image}{本质像}为 $D^+(S)$ 中由\Term{cohomology}{上同调}\Term{quasi-coherent}{拟凝聚}的\Term{complex}{复形}组成的\Term{full-subcategory}{全子范畴} $D^+(S)_{\mathrm{qcoh}}$。
\end{SGAresult}
\begin{proof}
只须证明 $S$ \Term{separated}{分离}的情形，因为 $S$ \Term{noetherian-scheme}{诺特}的情形已知（[H] II 7.19）。显然，这\Term{equivalence}{等价}于证明：

(3.5.1) 对 $E\in\operatorname{ob}D^+(\operatorname{Qcoh}(S))$，由 3.3 b)(ii) 得到的\Term{adjunction-morphism}{伴随箭头}
\[
E\longrightarrow\mathbf RQ(E)
\]
为\Term{isomorphism}{同构}；

(3.5.2) 对 $E\in\operatorname{ob}D^+(S)_{\mathrm{qcoh}}$，由 3.3 b)(ii) 得到的\Term{adjunction-morphism}{伴随箭头}
\[
\mathbf RQ(E)\longrightarrow E
\]
为\Term{isomorphism}{同构}。

\textbf{3.5.1 的证明。}直接逐层\Term{resolution}{分解}，\footnote{例如使用\Term{spectral-sequence}{谱序列} $R^pQ(H^q(E))\Longrightarrow R^*Q(E)$。}即归结到证明\Term{quasi-coherent}{拟凝聚}\Term{sheaf}{层}对\Term{coherator}{凝聚化函子}\Term{acyclic-sheaf}{无上同调}（即对 $E\in\operatorname{ob}\operatorname{Qcoh}(S)$，$R^iQ(E)=0$ 对 $i>0$ 成立），并且\Term{adjunction-morphism}{伴随箭头} $E\to QE$ 为\Term{isomorphism}{同构}。若 $S$ \Term{affine-scheme}{仿射}，这由公式 $Q=\widetilde{\Gamma_S}$ 得到（参见 3.2 的证明）。\Term{generic}{一般}地，以\Term{finite}{有限}个\Term{affine-scheme}{仿射}\Term{open-subset}{开集} $U_i$ \Term{covering}{覆盖} $S$。因 $S$ \Term{separated}{分离}，包含 $f_i:U_i\to S$ 均为\Term{affine-morphism}{仿射态射}，
\SourcePage{200}
各包含 $f_{i_0\cdots i_p}:U_{i_0\cdots i_p}=U_{i_0}\cap\cdots\cap U_{i_p}\to S$ 也一样。

设 $E\in\operatorname{ob}\operatorname{Qcoh}(S)$，考虑由\Term{covering}{覆盖} $(U_i)$ 定义的 $E$ 的\Term{resolution}{消解}（T.F. II 5.2）：
\[
0\longrightarrow E\longrightarrow\prod_i f_{i*}f_i^*(E)\longrightarrow\cdots\longrightarrow\prod_{i_0\cdots i_p}f_{i_0\cdots i_p*}f_{i_0\cdots i_p}^*(E)\longrightarrow\cdots.\tag{*}
\]
它给出\Term{spectral-sequence}{谱序列}
\[
E_1^{pq}=\prod_{i_0\cdots i_p}R^qQ\bigl(f_{i_0\cdots i_p*}f_{i_0\cdots i_p}^*(E)\bigr)\Longrightarrow R^*Q(E).\tag{**}
\]
因 $f_{i_0\cdots i_p}$ \Term{affine-morphism}{仿射}，由 3.4，对于每个 $q$ 有
\[
R^qQ\bigl(f_{i_0\cdots i_p*}f_{i_0\cdots i_p}^*(E)\bigr)
\simeq f_{i_0\cdots i_p*}R^qQ\bigl(f_{i_0\cdots i_p}^*(E)\bigr).\tag{***}
\]
特别地，由已处理的\Term{affine-scheme}{仿射}情形，$E_1^{pq}=0$ 对 $q>0$ 成立。因此由 (**) 得到\Term{isomorphism}{同构}
\[
R^nQ(E)\simeq H^n(QE^\bullet)\qquad(n\geq0),
\]
其中 $E^\bullet$ 为\Term{exact-sequence}{正合列} (*) 定义的 $E$ 的\Term{resolution}{消解}。但由 (***)，对 $n\geq0$ 有 $QE^n\simeq E^n$，最终得到
\[
R^nQ(E)=0\quad(n>0),\qquad Q(E)\simeq E,
\]
证毕。

\textbf{3.5.2 的证明。}由\Term{spectral-sequence}{谱序列}
\[
R^pQ(H^q(E))\Longrightarrow R^*Q(E),
\]
\SourcePage{201}
立即归结到 $E\in\operatorname{ob}\operatorname{Qcoh}(S)$ 的情形。\Term{adjunction-morphism}{伴随箭头}的复合
\[
E\longrightarrow\mathbf RQE\longrightarrow E
\]
为恒等，故断言由 3.5.1 得到。这完成了 3.5 的证明。
\end{proof}
\textbf{注记 3.6。}若去掉“\Term{separated}{分离}或\Term{noetherian-scheme}{诺特}”的假设，3.5 的结论不成立；J. L. Verdier 的一个漂亮反例说明了这一点（见本讲附录 I）。
\begin{SGAresult}{命题}{3.7}
在 3.5 的假设下，另假定存在整数 $N$，使对每个 $F\in\operatorname{ob}\operatorname{Mod}(S)$ 和 $S$ 的每个\Term{quasi-compact}{拟紧}\Term{open-subset}{开集} $U$，有
\[
H^i(U,F)=0\qquad(i>N).
\]
\footnote{由 Tohoku 3.6.5，若 $S$ 的底\Term{space}{空间}\Term{noetherian-space}{诺特}且 \Term{krull-dimension}{Krull 维数} $\leq N$，则此条件成立。}则：

a) \Term{coherator}{凝聚化函子}有\Term{finite}{有限}\Term{cohomological-dimension}{上同调维数}。

b) \Term{functor}{函子}
\[
\phi:D(\operatorname{Qcoh}(S))\longrightarrow D(S)
\]
\Term{fully-faithful}{全忠实}，其\Term{essential-image}{本质像}为 $D(S)$ 中由\Term{cohomology}{上同调}\Term{quasi-coherent}{拟凝聚}的\Term{complex}{复形}组成的\Term{full-subcategory}{全子范畴} $D(S)_{\mathrm{qcoh}}$。
\end{SGAresult}
\begin{proof}[a) 的证明]
若 $S$ \Term{affine-scheme}{仿射}，公式 $Q=\widetilde{\Gamma_S}$（见 3.2 的证明）表明 $Q$ 的\Term{cohomological-dimension}{上同调维数} $\leq N$。现考虑\Term{generic}{一般}情形。证明存在整数 $a$，使对所有 $E\in\operatorname{ob}\operatorname{Mod}(S)$ 和 $q>a$，有 $R^qQ(E)=0$。假设蕴含每个 $E\in\operatorname{ob}\operatorname{Mod}(S)$ 都有长为 $N$ 的\Term{finite}{有限}右\Term{resolution}{消解} $E^\bullet$，其各 $E^i$ \Term{acyclic-sheaf}{无上同调}，即对 $S$ 的每个\Term{open-subset}{开集} $U$ 及 $q>0$，均有 $H^q(U,E^i)=0$。故只须证明存在整数 $b$，使对每个\Term{acyclic-sheaf}{无上同调} $\mathcal O_S$-\Term{module}{模} $E$ 和 $q>b$，有 $R^qQ(E)=0$（这时可取 $a=b+N$）。取 $S$ 的一个由 $N'$ 个\Term{affine-scheme}{仿射}\Term{open-subset}{开集}组成的\Term{finite}{有限}\Term{covering}{覆盖} $(U_i)$。
\SourcePage{202}
设 $E$ 为\Term{acyclic-sheaf}{无上同调} $\mathcal O_S$-\Term{module}{模}，考虑由\Term{covering}{覆盖} $(U_i)$ 定义的 $E$ 的\Term{resolution}{消解} $C^\bullet(E)$（T.F. II 5.2）。分两种情形。

$1^\circ$）$S$ \Term{separated}{分离}。此时 $C^\bullet(E)$ 的各\Term{set-component}{分量}为形如 $f_*f^*E$ 的项的和，其中 $f:U\to S$ 为\Term{affine-scheme}{仿射}\Term{open-subset}{开集}的包含（因 $S$ \Term{separated}{分离}，它也是\Term{affine-morphism}{仿射态射}）。有
\[
\begin{aligned}
\mathbf R(Qf_*)f^*E&\simeq\mathbf RQ\,\mathbf Rf_*f^*E\\
&\simeq\mathbf RQ\,f_*f^*E,
\end{aligned}
\]
因为 $f^*E$ \Term{acyclic-sheaf}{无上同调}。另一方面，由 3.4，
\[
\begin{aligned}
\mathbf R(Qf_*)f^*E&\simeq\mathbf R(f_*Q)f^*E\\
&\simeq f_*\mathbf RQ\,f^*E,
\end{aligned}
\]
因为 $f$ \Term{affine-morphism}{仿射}。由此得
\[
R^nQ(f_*f^*E)\simeq f_*R^nQ(f^*E)\qquad(n\geq0),
\]
进而由\Term{affine-scheme}{仿射}情形，
\[
R^nQ(f_*f^*E)=0\qquad(n>N).
\]
因 $C^\bullet(E)$ 长为 $N'$，最终有
\[
R^nQ(E)=0\qquad(n>N+N'),
\]
结论得证。

$2^\circ$）$S$ \Term{noetherian-scheme}{诺特}。$C^\bullet(E)$ 的各\Term{set-component}{分量}为形如 $f_*f^*E$ 的项的和，其中 $f:U\to S$ 为\Term{quasi-compact}{拟紧}\Term{separated}{分离}\Term{open-subset}{开集}的包含。与情形 $1^\circ$ 相同，有
\SourcePage{203}
\[
\mathbf R(Qf_*)f^*E\simeq\mathbf RQ\,f_*f^*E\simeq\mathbf R(f_*Q)f^*E.\tag{*}
\]
因 $U$ \Term{noetherian-scheme}{诺特}，$\operatorname{Qcoh}(U)$ 有足够多在 $\operatorname{Mod}(U)$ 中\Term{injective-object}{内射}的\Term{object}{对象}（[H] II 7.18）；故经由 3.5 将 $D^+(\operatorname{Qcoh}(U))$（相应地 $D^+(\operatorname{Qcoh}(S))$）等同于 $D^+(U)$（相应地 $D^+(S)$）的\Term{full-subcategory}{全子范畴}时，$f_*:\operatorname{Qcoh}(U)\to\operatorname{Qcoh}(S)$ 的\Term{right-derived-functor}{右导出函子}由 $\mathbf Rf_*:D^+(U)\to D^+(S)$ “诱导”。所以有
\[
\mathbf R(f_*Q)f^*E\simeq\mathbf Rf_*\,\mathbf RQ\,f^*E.
\]
从而得到\Term{spectral-sequence}{谱序列}
\[
E_2^{pq}=R^pf_*\,R^qQ\,f^*E\Longrightarrow R^*(f_*Q)f^*E.
\]
由情形 $1^\circ$，存在仅依赖于 $U$ 的整数 $n(U)$，使 $Q$ 在 $\operatorname{Mod}(U)$ 上的\Term{cohomological-dimension}{上同调维数} $\leq n(U)$。于是
\[
E_2^{pq}=0\qquad(q>n(U)).
\]
但由关于 $S$ 的“\Term{cohomological-dimension}{上同调维数}”的假设，还有
\[
E_2^{pq}=0\qquad(p>N).
\]
因此
\[
R^n(f_*Q)f^*E=0\qquad(n>N+N''),
\]
其中 $N''$ 为大于所有 $n(U)$ 的整数，$U$ 遍历各交集 $U_{i_0}\cap\cdots\cap U_{i_p}$。结合 (*)，得到
\[
R^nQ(f_*f^*E)=0\qquad(n>N+N''),
\]
最后得到
\[
R^nQ(E)=0\qquad(n>N+N'+N''),
\]
\SourcePage{204}
这完成了 a) 的证明。
\end{proof}
\begin{proof}[b) 的证明]
由 a)，\Term{coherator}{凝聚化函子}有\Term{finite}{有限}\Term{cohomological-dimension}{上同调维数}，故有\Term{right-derived-functor}{右导出函子}（[H] I 5.3）
\[
\mathbf RQ:D(S)\longrightarrow D(\operatorname{Qcoh}(S)),
\]
其在 $D^+(S)$ 上的\Term{restriction}{限制}与上述 $\mathbf RQ$ 一致。对 $E\in\operatorname{ob}D(\operatorname{Qcoh}(S))$，有\Term{canonical}{典范}\Term{functorial}{函子性}\Term{isomorphism}{同构}
\[
E\xrightarrow{\sim}\mathbf RQ(E),\tag{1}
\]
因为由 3.5.1，$E^i$ 对 $Q$ \Term{acyclic-sheaf}{无上同调}，且 $E^i\xrightarrow{\sim}Q(E^i)$。另一方面，对 $E\in\operatorname{ob}D(S)$，有\Term{canonical}{典范}\Term{functorial}{函子性}\Term{homomorphism}{同态}
\[
\mathbf RQ(E)\longrightarrow E,\tag{2}
\]
当 $E^i$ 对 $Q$ \Term{acyclic-sheaf}{无上同调}时，它来自 $\phi$ 与 $Q$ 之间通常的\Term{adjunction-morphism}{伴随同态}。容易验证，(1)、(2) 使 $\mathbf RQ:D(S)\to D(\operatorname{Qcoh}(S))$ 成为 $\phi:D(\operatorname{Qcoh}(S))\to D(S)$ 的\Term{right-adjoint}{右伴随}。(1) 为\Term{isomorphism}{同构}这一事实表明 $\phi$ \Term{fully-faithful}{全忠实}。最后，使用\Term{spectral-sequence}{谱序列}
\[
R^pQ(H^q(E))\Longrightarrow R^*Q(E)
\]
（它因 $Q$ 有\Term{finite}{有限}\Term{cohomological-dimension}{上同调维数}而收敛），可知对 $E\in\operatorname{ob}D(S)_{\mathrm{qcoh}}$，(2) 为\Term{isomorphism}{同构}；因此 $\phi$ 的\Term{essential-image}{本质像}为 $D(S)_{\mathrm{qcoh}}$。这证明了 b)，并完成了 3.7 的证明。
\end{proof}
